%==============================================================================
% Dynamic Career Signaling, Employer Learning, and Optimal Labor Income Taxation
% Evidence from the NLSY79
% Draft — NLSY79 Replication and Extension of Sztutman (2024)
%==============================================================================

\documentclass[12pt]{article}

\usepackage[margin=1.25in]{geometry}
\usepackage{setspace}
\usepackage{amsmath, amssymb, amsthm}
\usepackage{graphicx}
\usepackage{booktabs}
\usepackage{threeparttable}
\usepackage{caption}
\usepackage{subcaption}
\usepackage{natbib}
\usepackage{hyperref}
\usepackage{xcolor}
\usepackage{microtype}
\usepackage{lmodern}
\usepackage{float}
\usepackage{array}
\usepackage{dcolumn}

\doublespacing

\hypersetup{
    colorlinks=true,
    linkcolor=black,
    citecolor=black,
    urlcolor=blue
}

\newtheorem{proposition}{Proposition}
\newtheorem{corollary}{Corollary}[proposition]
\theoremstyle{definition}
\newtheorem{assumption}{Assumption}
\theoremstyle{remark}
\newtheorem*{remark}{Remark}

% Possessive citation: \citepos{key} → Author's (Year)
\newcommand{\citepos}[1]{\citeauthor{#1}'s (\citeyear{#1})}

% Column type for decimal-aligned numbers
\newcolumntype{d}[1]{D{.}{.}{#1}}

%==============================================================================
\begin{document}
%==============================================================================

\begin{titlepage}
\vspace*{2cm}
\begin{center}
{\Large \textbf{Dynamic Career Signaling, Employer Learning, and Optimal Labor Income Taxation: Evidence from the NLSY79}}

\vspace{1.5cm}

{\large Yuhao Ren}\\
{\normalsize Carnegie Mellon University}\\
{\normalsize Tepper School of Business}

\vspace{0.5cm}
{\normalsize Advised by Prof.\ Andr\'{e} Sztutman}

\vspace{1.5cm}
{\normalsize \textit{Preliminary Draft — Do Not Circulate}}\\
{\normalsize April 2026}

\vspace{2cm}
\\[3em]
\begin{minipage}{0.85\textwidth}
\small\textbf{Abstract.}
Between 1979 and 2023, the wages of NLSY79 workers grew at an elasticity of 0.94 with
respect to their cumulative hours worked---a return that declines from 0.97 at career entry
to 0.83 at midcareer. This pattern is the empirical fingerprint of employer learning: as
workers accumulate a public labor market history, employers gradually resolve their uncertainty
about ability, and the informational value of each additional hour diminishes.

Employer learning creates a positive externality that standard markets do not price. Each
hour a worker supplies is recorded in her employment history and observed by all future
employers at no cost---but only the current employer compensates her for it. Because workers
capture only part of the social value of the information they generate, labor is undersupplied
relative to the social optimum. The standard remedy is a Pigouvian subsidy to labor income.
Sztutman (2024) shows that this subsidy takes the closed-form $\tau_p = \delta / \alpha$,
where $\delta$ measures the fraction of wage growth driven by ongoing employer learning rather
than genuine skill accumulation, and $\alpha$ is the Pareto tail index of the income
distribution. The tail index enters because the Pareto distribution's constant-ratio property
pins down the externality as a fixed fraction $1/\alpha$ of wages at every income level.

I estimate both objects from NLSY79 panel data covering 12,686 workers, using person
fixed-effects wage regressions to recover $\delta$ and federal tax reform
variation---spanning five reforms from ERTA (1981) through TCJA (2017)---as instruments for
marginal tax rate changes to recover workers' behavioral response to taxation. The primary
estimate is a Pigouvian subsidy of 19.1 percent of labor income (95\% CI: [15.6\%, 22.6\%]),
substantially larger than prior estimates for older workers, consistent with the prediction
that employer uncertainty is greatest early in careers. Occupation-level variation in
wage-hours elasticities and the finding that the wage premium for pre-labor-market ability
rises with career experience---a direct implication of gradual employer learning---corroborate
the mechanism.
\end{minipage}
\vspace{2cm}\\
\small\textit{Keywords: employer learning, signaling, optimal taxation, elasticity of taxable income, Pigouvian tax}
\end{center}
\end{titlepage}

%==============================================================================
\section{Introduction}
\label{sec:intro}
%==============================================================================

Between 1979 and 1993, the median real wage of young American workers nearly doubled, growing
at a rate that substantially exceeds what accumulated human capital can explain. In the National
Longitudinal Survey of Youth 1979 cohort (NLSY79), each 10 percent increase in cumulative hours
worked is associated with a 9.4 percent increase in wages after controlling for potential
experience, year fixed effects, and unobserved worker ability through person fixed effects. This
elasticity—call it $\gamma$—is not constant over careers: it equals 0.97 at labor market entry,
declines monotonically to a minimum of 0.83 around 22 years of experience, then rises slightly
at the end of the sample window. The systematic mid-career decline cannot be explained by
diminishing returns to human capital accumulation, which would require the observed elasticity
to be stable or smoothly declining throughout. It is, however, exactly what dynamic signaling
models predict: as workers accumulate a public work history, employers learn progressively more
about their underlying ability, and the incremental informational value of each additional hour
of work diminishes.

The employer learning interpretation has direct implications for optimal labor income taxation
that have not been quantified. When a worker increases labor supply, she generates information
about her ability that is partly observed by all future employers, not only by the current
employer who compensates her for the additional work. This positive externality drives a wedge
between the private return to labor effort and its social return, creating a standard
Pigouvian case for a labor income subsidy \citep{boskin1975taxation, stiglitz1982self}.
\citet{sztutman2024} formalizes this intuition in a dynamic signaling model with Pareto-distributed
productivity types and derives a closed-form expression for the optimal Pigouvian subsidy:
$\tau_p = \delta / \alpha$, where $\delta \in (0,1)$ measures the fraction of wage growth
driven by ongoing employer learning rather than genuine skill accumulation (higher $\delta$
means more residual employer uncertainty), and $\alpha$ is the Pareto tail index of the income
distribution. The tail index enters because the Pareto distribution has a distinctive
constant-ratio property---the expected income above any threshold is proportional to the
threshold itself, with constant $\alpha/(\alpha-1)$---which pins down the externality magnitude
as exactly $1/\alpha$ of wages at every income level. The key insight of this framework is
that the optimal subsidy depends on the data only through the pair $(\delta, \alpha)$, and
$\delta$ is in turn identified by two observables: the wage-hours elasticity $\gamma$
(recoverable from panel wage regressions) and the elasticity of taxable income $\varepsilon^w$
(recoverable from tax reform variation). No direct measure of employer beliefs or worker types
is required. Yet no study has estimated this formula empirically on a representative cohort of
workers over their full careers.

This paper fills that gap. I make four contributions. First, I estimate the career arc of the
wage-cumulative-hours elasticity $\gamma$ using a polynomial fixed-effects specification that
recovers a continuous $\gamma(\text{exp})$ function without imposing the discretization error
of quintile bins. The estimated arc—declining from 0.97 at entry to 0.83 at midcareer—is
inconsistent with the Pareto constant-$\gamma$ case (rejected at $F(2, 11524) = 43.20$) and
consistent with the dynamic employer learning prediction.
Second, I estimate how workers' reported earnings respond to changes in their marginal
tax rate---the elasticity of taxable income---in two distinct career phases, using changes
in federal and state tax schedules driven by five major legislative reforms as instruments
\citep{gruber2002elasticity}. The annual period (1978--1993, ages 17--35) yields a
statistically insignificant elasticity of $-0.048$ (SE $= 0.080$); the biennial period
(1995--2019, ages 31--62) yields a significant elasticity of $+0.297$ (SE $= 0.121$, $t =
2.46$). The contrast reflects a lifecycle pattern: young workers optimize reputation
accumulation rather than tax liability, with the behavioral response to taxes emerging only
as careers mature and employer uncertainty resolves.
Third, I combine these sufficient statistics with a directly estimated Pareto tail parameter
$\hat{\alpha} = 3.295$ (SE $= 0.014$, $R^2 = 0.985$) to produce the first structural
panel-data estimate of the Pigouvian labor income subsidy identified from lifecycle variation
in both $\hat{\gamma}$ and $\hat{\varepsilon}^w$: $\hat{\tau}_p = 19.1\%$ (95\% CI:
$[15.6\%$, $22.6\%]$). Unlike \citepos{sztutman2024} illustrative calculation for late-career
HRS respondents, which calibrates the sufficient statistics to approximate values rather than
estimating them from an identified panel regression design, this estimate exploits federal tax
reform variation to instrument for marginal rate changes and recovers the structural $\delta$
from a representative birth cohort observed over their full careers.
Fourth, I cross-validate
the employer learning mechanism through sector heterogeneity in $\gamma$—which increases
monotonically from 0.90 in low-signaling occupations to 1.04 in high-signaling occupations—and
through the Altonji--Pierret (2001) test for gradual employer learning---which predicts
that the wage premium for pre-labor-market ability, as measured by Armed Forces Qualifying
Test scores, should rise with career experience as the market progressively incorporates
information about worker type---finding exactly this pattern among workers with established
occupation tenure. Appendix Table~\ref{tab:gamma_occ} reports $\hat{\gamma}$ for all nine 1-digit
occupations; Sales workers have the highest $\hat{\gamma} = 1.235$ (output observability
problems in client-facing roles) and Farm Workers the lowest $\hat{\gamma} = 0.470$ (directly
observable agricultural output), anchoring the group-level pattern at its extremes.

The magnitude of the estimated subsidy deserves emphasis. At 19.1 percent of labor income, the
implied Pigouvian correction is economically consequential: it implies that the optimal marginal
labor income tax rate is 19 percentage points below the rate dictated by equity and revenue
considerations alone. This estimate is substantially larger than \citepos{sztutman2024}
back-of-envelope estimate of approximately 5 percent for workers ages 50 and above in the
Health and Retirement Study (HRS). The gap is predicted by the model: the signaling
externality is strongest early in careers, when employers face the most uncertainty, and
declines as careers mature and reputations are established. The implied signaling parameter
$\hat{\delta} = 0.628$ for NLSY79 workers (average age 33) versus Sztutman's $\delta \approx 0.17$
for HRS workers (average age 60) traces a lifecycle arc in which 73 percent of employer
uncertainty resolves between early and late career. Cross-cohort extrapolation suggests that
the average Pigouvian subsidy, integrated over a full career, lies in the range of 10--15
percent of labor income.

The identification strategy relies on three separate sources of variation. The structural
$\gamma$ is identified from within-person variation in the wage-cumulative-hours relationship,
purged of unobserved worker heterogeneity through person fixed effects (N $= 179{,}209$
person-years). The elasticity of taxable income is identified from simulated marginal tax rate
changes driven by federal tax reforms---ERTA (1981), TRA (1986), EGTRRA (2001), JGTRRA (2003),
and TCJA (2017)---following \citet{gruber2002elasticity}, with first-stage $F$-statistics of
2,617 (annual) and 1,146 (biennial), well above conventional weak-instrument thresholds. The
Pareto tail is estimated via log-rank regression on the top 20 percent of the annual income
distribution. All three estimation steps use clustered standard errors at the person level.

This paper contributes to four literatures. First, it extends the employer learning literature
\citep{farber1996learning, altonji2001employer, gibbons1999careers, lange2007speed} by providing
the first estimate of the Pigouvian tax implied by the learning process, linking the
well-documented empirical phenomenon to its optimal tax consequence. Second, it extends the
optimal taxation of signaling literature \citep{boskin1975taxation, stiglitz1982self,
sztutman2024} by supplying the missing empirical inputs to the tax formula; prior work derived
the formula theoretically but left quantification to future research. Third, it adds to the
elasticity of taxable income literature \citep{feldstein1995effect, gruber2002elasticity,
kopczuk2005tax, saez2012elasticity} by documenting a lifecycle regime change in the ETI—near
zero for young workers, positive for prime-career workers—that reflects the signaling mechanism
rather than a single behavioral elasticity. Relative to this literature, the novel finding is
not the magnitude of the ETI but the lifecycle pattern and its structural interpretation.
Fourth, it provides the first direct estimate of $\alpha$ from NLSY79 earnings data, adding to
the empirical literature on Pareto tail estimation \citep{saez2001using, diamond1998optimal}.

The paper proceeds as follows. Section~\ref{sec:model} presents the Sztutman (2024) model and
derives the sufficient-statistic formula for $\tau_p$. Section~\ref{sec:data} describes the
NLSY79 data and sample construction. Section~\ref{sec:empirics} specifies the three estimation
equations and describes the identification strategy for each. Section~\ref{sec:results} reports
the main estimates: the $\gamma$ career arc, the two-period ETI, and the implied $\tau_p$.
Section~\ref{sec:validation} provides cross-validation evidence through sector heterogeneity
and the AP test. Section~\ref{sec:robustness} conducts robustness checks addressing
collinearity, simultaneity, heterogeneous trends, and sensitivity to the Pareto assumption.
Section~\ref{sec:discussion} discusses the lifecycle arc of $\delta$, the comparison to
Sztutman (2024), and the policy implications. Section~\ref{sec:conclusion} concludes.

%==============================================================================
\section{Theoretical Framework}
\label{sec:model}
%==============================================================================

\subsection{Setup}

I adopt the dynamic signaling framework of \citet{sztutman2024}, which builds on the
job-market signaling foundation of \citet{spence1973job} by introducing dynamic employer
learning and an explicit role for cumulative work history as a signal. A worker of type
$\theta$ enters the labor market at time 0. The type distribution is:

\begin{assumption}
\label{ass:pareto}
Worker types are drawn from a Pareto distribution: $1 - F(\theta) = (\underline{\theta}/\theta)^\alpha$
for $\theta \geq \underline{\theta} > 0$, where $\alpha > 1$ is the tail parameter and
$\underline{\theta}$ is the minimum type. Type $\theta$ is private information to the worker.
\end{assumption}

Employers observe only the worker's public labor market history. At each period $t$, the
worker chooses hours $h_t \geq 0$; cumulative hours through period $t$ are
$H_t = \sum_{s \leq t} h_s$. Employers update beliefs about worker type by observing $H_t$:

\begin{assumption}
\label{ass:learning}
Beliefs update monotonically in work history: $E[\theta \mid H_t]$ is strictly increasing in
$H_t$ for all $t$. Wages equal the employer's expected marginal product:
$w_t = E[f(\theta) \mid H_t]$, where $f(\cdot)$ is the worker's production function.
\end{assumption}

Assumption~\ref{ass:learning} is satisfied under Bayesian learning from a Gaussian signal
structure \citep{farber1996learning} as well as under the specific signal-extraction structure
in \citet{sztutman2024}. It implies that accumulating a longer public work history increases
wages by signaling higher type---the core mechanism of the model.

Under Assumption~\ref{ass:learning}, log-linearizing around the steady-state wage schedule
yields the estimating equation:

\begin{equation}
\label{eq:wage}
\log w_{it} = \alpha_i + \gamma \cdot \log H_{it} + \beta X_{it} + \varepsilon_{it},
\end{equation}

where $\alpha_i$ is a worker fixed effect capturing time-invariant ability as priced by the
market at career entry, $\gamma$ is the wage-cumulative-hours elasticity, $X_{it}$ includes
potential experience, its square, and year fixed effects, and $\varepsilon_{it}$ is mean-zero
idiosyncratic error. The coefficient $\gamma$ is the central object of interest: it maps the
informational content of the public work history into within-person wage growth. A worker who
doubles her cumulative hours raises her wage by $\gamma$ percent, holding within-person ability
pricing fixed.

\subsection{Structural Interpretation of $\gamma$}

Under Assumptions~\ref{ass:pareto}--\ref{ass:learning}, \citet{sztutman2024} derives the
structural link between the observable elasticity $\gamma$ and the model's primitive
parameters. Let $\delta \in (0,1)$ denote the probability that employers remain uninformed
about a worker's type after a given period---the \emph{signaling parameter} measuring the
residual employer uncertainty. \citet{sztutman2024} shows:

\begin{equation}
\label{eq:gamma_structural}
\gamma = \frac{\delta}{1 - \delta + \varepsilon^w},
\end{equation}

where $\varepsilon^w$ is the elasticity of taxable income. Equation~(\ref{eq:gamma_structural})
has a direct intuition. The numerator $\delta$ is the share of the wage return to hours that
flows through the signaling channel: higher $\delta$ means employers update beliefs more
aggressively in response to additional work history, raising $\gamma$. The denominator
$1 - \delta + \varepsilon^w$ normalizes by the total margin of adjustment: as $\delta
\to 0$ (complete employer knowledge), $\gamma \to 0$; as $\delta \to 1$ (complete employer
ignorance), $\gamma \to 1/\varepsilon^w$.

Inverting equation~(\ref{eq:gamma_structural}) gives:

\begin{equation}
\label{eq:delta}
\delta = \frac{\gamma(1 + \varepsilon^w)}{1 + \gamma}.
\end{equation}

Equation~(\ref{eq:delta}) shows that $\delta$ is identified by two observables: the
structural elasticity $\gamma$, recoverable from the panel wage regression (\ref{eq:wage}),
and the ETI $\varepsilon^w$, recoverable from tax reform variation. Both are estimable from
standard panel data without requiring direct measures of employer beliefs.

\subsection{The Pigouvian Tax Formula}

The social planner maximizes a utilitarian welfare function subject to the private optimality
conditions of workers. The planner recognizes what individual workers do not internalize: each
additional hour of work generates information about ability that is observed by all future
employers, creating a positive externality. The private return to an additional hour equals
the current wage; the social return includes the current wage \emph{plus} the discounted value
of improved employer beliefs across all future employment relationships. This wedge between
private and social returns is the Pigouvian externality.

\citet{sztutman2024} shows that under Assumptions~\ref{ass:pareto}--\ref{ass:learning},
the Pareto income distribution pins down the magnitude of this externality in closed form.
The optimal Pigouvian tax formula is:

\begin{equation}
\label{eq:taup}
\tau_p = \frac{\delta}{\alpha}.
\end{equation}

\begin{proposition}[\citealt{sztutman2024}, Proposition~3]
\label{prop:taup}
Under Assumptions~\ref{ass:pareto}--\ref{ass:learning}, with a benevolent social planner
who can observe work history but not worker type, the optimal labor income tax rate is
$t^* = t_{\text{equity}} - \tau_p$, where $t_{\text{equity}}$ is the rate dictated by
redistributive considerations alone and
\[
\tau_p = \frac{\delta}{\alpha}.
\]
The optimal subsidy $\tau_p$ is (i) increasing in $\delta$: a higher probability of residual
employer ignorance generates a larger externality; (ii) decreasing in $\alpha$: a heavier
Pareto tail (smaller $\alpha$) means that the income gains from improved employer beliefs
are concentrated among fewer high-type workers, reducing the average externality; and
(iii) independent of the income floor $\underline{\theta}$ and the functional form of $f$.
\end{proposition}

\noindent \emph{Proof sketch.} The planner's problem reduces to choosing a labor income tax
$t$ to maximize worker welfare subject to participation constraints. The first-order condition
generates a standard Ramsey rule augmented by the Pigouvian externality term. Under the Pareto
distribution (Assumption~\ref{ass:pareto}), the externality integral evaluates to $\delta / \alpha$
because the Pareto tail implies a constant ratio of the marginal worker's type to the average
type above her. This constant-ratio property---characteristic of the Pareto and no other
distribution---is what makes the formula scale-free. Full proof: see \citet{sztutman2024},
Section~3.2 and Appendix~B.

Proposition~\ref{prop:taup} is a \emph{sufficient statistic} result in the spirit of
\citet{chetty2009sufficient}: the welfare-optimal tax adjustment depends on the data only
through the pair $(\delta, \alpha)$, and $\delta$ is in turn identified by the observable
pair $(\gamma, \varepsilon^w)$ via equation~(\ref{eq:delta}). The structural parameters
$(\gamma, \varepsilon^w, \alpha)$ are therefore \emph{all that is needed} to quantify the
optimal labor income subsidy implied by employer learning---no direct measure of employer
beliefs or worker types is required. Section~\ref{sec:empirics} develops the identification
strategy for each of these three objects.

\subsection{Career Dynamics of $\gamma$ and an Empirical Prediction}

Equation~(\ref{eq:gamma_structural}) immediately implies that $\gamma$ should vary over the
career if $\delta$ varies with experience. Under the dynamic signaling interpretation, employer
uncertainty $\delta(\text{exp})$ declines as careers progress: at labor market entry, employers
have no work history to observe and face maximal uncertainty; as the worker accumulates a
longer public record, beliefs converge toward the true type and $\delta$ falls. Holding
$\varepsilon^w$ approximately constant, equation~(\ref{eq:gamma_structural}) delivers a
testable prediction:

\begin{corollary}[Dynamic signaling prediction]
\label{prop:gamma_arc}
Under Assumptions~\ref{ass:pareto}--\ref{ass:learning}, if $\delta(\text{exp})$ is strictly
decreasing in potential experience, then $\gamma(\text{exp}) = \delta(\text{exp}) /
(1 - \delta(\text{exp}) + \varepsilon^w)$ is also strictly decreasing in experience.
In the limit as $\delta(\text{exp}) \to 0$ (complete resolution of employer uncertainty),
$\gamma(\text{exp}) \to 0$.
\end{corollary}

\noindent \emph{Derivation.} Differentiating equation~(\ref{eq:gamma_structural}) with
respect to experience:
\[
\frac{\partial \gamma}{\partial \text{exp}}
= \frac{(1 - \delta + \varepsilon^w) \delta'(\text{exp}) - \delta \cdot (-\delta'(\text{exp}))}{(1 - \delta + \varepsilon^w)^2}
= \frac{\delta'(\text{exp})}{(1 - \delta + \varepsilon^w)^2} \cdot (1 + \varepsilon^w).
\]
Since $1 + \varepsilon^w > 0$ (the tax base is not Laffer-backward) and $(1 - \delta + \varepsilon^w)^2 > 0$,
the sign of $\partial \gamma / \partial \text{exp}$ equals the sign of $\delta'(\text{exp})$.
If $\delta'(\text{exp}) < 0$, then $\gamma'(\text{exp}) < 0$. $\square$

This prediction distinguishes the signaling model from pure human capital theory. Under
human capital theory, $\gamma$ reflects the marginal return to skill accumulation; absent
a model of depreciation or diminishing returns to on-the-job learning, $\gamma$ has no
systematic career arc once person fixed effects absorb time-invariant ability. The
corollary therefore provides a falsifiable joint test of the employer learning mechanism:
if the estimated $\hat{\gamma}(\text{exp})$ is monotone-declining, employer learning is
the dominant process driving wage-hours covariation over the career. If $\hat{\gamma}(\text{exp})$
is flat or increasing, either employer learning is absent or some other mechanism---such as
on-the-job training accumulation or match quality deterioration---offsets the informational
convergence.

%==============================================================================
\section{Data}
\label{sec:data}
%==============================================================================

\subsection{The NLSY79 Panel}

The primary data source is the National Longitudinal Survey of Youth 1979 (NLSY79), a
nationally representative probability sample of 12,686 individuals born between 1957 and 1964.
The Bureau of Labor Statistics designed the NLSY79 to follow respondents from labor market
entry through retirement; the survey collected annual observations from 1979 through 1994
and has continued biennially from 1996 onward. I use all survey rounds through 2023, yielding
up to 33 observations per person.

Two features make the NLSY79 well-suited for this study. First, the panel spans the full
career arc from labor market entry (respondents were 14--22 years old in 1979) through late
career (ages 58--65 in 2023), enabling estimation of $\hat{\gamma}(\text{exp})$ from
Corollary~\ref{prop:gamma_arc} across the complete experience range. Second, the respondents
were young adults during the major tax reforms of the 1980s---specifically, ERTA (1981)
cut the top marginal federal rate from 70 to 50 percent, and TRA (1986) cut it further to
28 percent---and were in their prime earning years during the major reforms of the 2000s
and 2010s. This lifecycle alignment provides multiple reform episodes that identify the ETI
at different career stages.

\paragraph{Attrition.} The NLSY79 experiences nontrivial attrition over 45 years of
follow-up. Of the original 12,686 respondents, roughly 69 percent have at least one
observation after 1995. Respondents who attrit earlier tend to have lower wages and education
than those who remain---a pattern of positive survival selection. This attrition is unlikely
to affect the structural $\gamma$ estimates, which exploit within-person wage-hours covariation
rather than cross-sectional levels. For the biennial ETI estimation (1995--2019), attrition
could cause sample selection bias if exits are correlated with the identifying tax variation;
however, the simulated instrument is constructed from the tax schedule at prior-period income,
not current-period income, so it is orthogonal to the selection decision at the time of exit.

\subsection{Key Variables}

\paragraph{Primary wages.} The wage variable $w_{it}$ is self-reported annual wage and salary
income from the primary job, deflated to 1984 constant dollars using the CPI-U. This is the
pre-tax income measure. I exclude person-years with non-positive wages from the structural
estimation; for the ETI estimation, the taxable income measure $z_{it}$ is wages minus the
applicable standard deduction (computed by TAXSIM35 given filing status).

\paragraph{Potential experience and education.} Potential experience is defined as
$\text{pot\_exp}_{it} = \text{age}_{it} - \text{hgc}_{it} - 6$, where $\text{hgc}_{it}$
is the highest grade completed as of the survey date (time-varying in the NLSY79). This
standard formula assigns 6 years to pre-school. I use the time-varying education measure
rather than the maximum-ever value to avoid attributing post-entry schooling to market
experience; the difference is small because the vast majority of NLSY79 respondents
complete their education before age 25.

\paragraph{Cumulative hours.} The cumulative hours variable $H_{it}$ is the sum of reported
annual hours worked from labor market entry through period $t$: $H_{it} = \sum_{s=1}^{t} h_{is}$.
Annual hours are directly reported in the NLSY79 for 1979--1994. For the biennial rounds
(1996 onward), I interpolate a linear trend between adjacent observed annual hours values
within each person, then accumulate. This linear interpolation is conservative: it assumes
no change in hours patterns between survey rounds and does not extrapolate beyond the
observed career window.

A specific correction addresses person-years with positive reported wages but zero or missing
reported hours---likely employed respondents who did not complete the hours module. For these
observations (approximately 3 percent of the annual-period sample), I impute hours using the
same person's median hours from valid adjacent observations. This correction raises early-career
cumulative hours slightly and moderates early-career $\hat{\gamma}$ estimates relative to a
specification that drops these observations.

\paragraph{Marginal tax rates via TAXSIM35.} Federal and state marginal income tax rates
are computed using TAXSIM35 \citep{feenberg1993introduction}, the standard NBER tax
simulation model used in the ETI literature \citep{gruber2002elasticity, kopczuk2005tax}.
For each person-year, I submit to TAXSIM35 the following inputs: filing status (single or
married filing jointly, based on NLSY79 marital status codes mapped to TAXSIM conventions);
age of primary taxpayer and spouse; number of dependents; primary wages; spousal wages
(where applicable); and state of residence. TAXSIM35 returns the marginal federal rate,
state rate, and total combined rate; I use the combined marginal rate throughout.

TAXSIM35 requires valid input ranges for all variables. Prior to submission, I apply the
following validation: (i) spouse age set to zero for non-married filers; (ii) dependent
count capped at 15; (iii) all income variables set to zero if negative or missing;
(iv) filing status forced to 1 (single) if the value is neither 1 nor 2. These constraints
follow the TAXSIM35 documentation and affect fewer than 2 percent of observations.

I construct the simulated instrument by resubmitting each observation to TAXSIM35 with
income held fixed at its prior-period value, inflated by the national average wage growth
rate, while all demographic characteristics are updated to their current-period values
(following \citealt{gruber2002elasticity}). This procedure generates a marginal rate that
reflects only legislative changes in the tax schedule, not changes in individual income.

\subsection{Sample Construction}

\paragraph{Structural sample.} The $\gamma$ estimation sample requires positive primary wages,
age 18--65, and potential experience in $[0, 40]$ years. This yields 179,209 person-years
across 12,686 individuals.

\paragraph{Annual ETI sample.} The ETI estimation requires pairs of consecutive annual
observations. I additionally apply: (i) an income floor of \$5,000 (1984 dollars) in both
the base and end year, to exclude workers whose reported income is likely below the TAXSIM35
valid range or dominated by non-wage income; (ii) a trimming rule $|\Delta \log z_{it}| \leq
\log 5$ to remove extreme income changes that likely reflect reporting errors, major
life events (divorce, disability), or transitory shocks rather than behavioral responses to
tax changes. These restrictions follow standard practice in the ETI literature \citep{gruber2002elasticity,
saez2012elasticity} and yield 53,632 person-years. Robustness to a lower \$1,000 floor is
reported in Section~\ref{sec:results}.

\paragraph{Biennial ETI sample.} The biennial estimation pairs observations separated by one
biennial survey interval (approximately 2 years). I apply a \$10,000 floor---consistent with
the higher absolute income levels for prime-career workers---and the same trimming rule,
yielding 47,201 person-years spanning 1995--2019.

Table~\ref{tab:desc} reports summary statistics for the full analytical population
(N $= 189{,}671$ person-years). The mean annual wage is \$27,695 (median \$17,000), reflecting
the right skew of the earnings distribution. Mean potential experience is 14.0 years and
mean cumulative hours are 26,344 hours, equivalent to roughly 13 years of full-time work.

%==============================================================================
\section{Empirical Strategy}
\label{sec:empirics}
%==============================================================================

\subsection{Estimating $\gamma$: Within-Person Fixed Effects}

The primary equation for $\gamma$ is:
\begin{equation}
\label{eq:fe_main}
\log w_{it} = \alpha_i + \gamma \cdot \log H_{it} + \beta_1 \cdot \text{pot\_exp}_{it}
+ \beta_2 \cdot \text{pot\_exp}_{it}^2 + \sum_t \lambda_t + \varepsilon_{it},
\end{equation}
estimated by within-person (FE) transformation on the structural sample. The person fixed
effect $\alpha_i$ absorbs all time-invariant sources of worker heterogeneity, including
unobserved ability. Standard errors are clustered at the person level throughout.

To estimate the career arc of $\gamma$, I extend equation~(\ref{eq:fe_main}) by interacting
$\log H_{it}$ with a second-order polynomial in potential experience:
\begin{equation}
\label{eq:fe_poly}
\log w_{it} = \alpha_i + \underbrace{(\beta_0 + \beta_1 \cdot \text{exp} + \beta_2 \cdot \text{exp}^2)}_{\gamma(\text{exp})} \cdot \log H_{it}
+ \text{controls} + \varepsilon_{it}.
\end{equation}
The implied elasticity at experience level $e$ is $\hat{\gamma}(e) = \hat{\beta}_0 +
\hat{\beta}_1 e + \hat{\beta}_2 e^2$, and the career minimum is $e^* = -\hat{\beta}_1 /
(2\hat{\beta}_2)$. A joint test of $H_0\colon \beta_1 = \beta_2 = 0$ tests whether $\gamma$
is constant over the career (the Pareto case).

\paragraph{Identification.} Within-person variation in $\log H_{it}$ identifies $\gamma$
after purging the fixed effect. The primary concern is collinearity: $\log H_{it}$ and
$\text{pot\_exp}_{it}$ both grow over the career, creating within-person correlation. I
address this through Frisch--Waugh--Lovell (FWL) verification: I orthogonalize $\log H_{it}$
on all other regressors within each person and re-estimate $\gamma$ on the residuals. FWL
guarantees identical point estimates, confirming the baseline specification is not mechanically
constructing $\hat{\gamma}$ through collinearity (Section~\ref{sec:robustness}). I also
instrument $\log H_{it}$ with its one-year lag to address within-year simultaneity
(Section~\ref{sec:robustness}).

\subsection{Estimating $\varepsilon^w$: Gruber--Saez IV}

\paragraph{Specification.} I estimate the elasticity of taxable income using the
two-stage least squares specification of \citet{gruber2002elasticity}. The second stage is:
\begin{equation}
\label{eq:eti}
\Delta \log z_{it} = \varepsilon^w \cdot \widehat{\Delta \log(1 - \tau_{it})} + \beta \cdot \log z_{i,t-k} + \alpha_t + \nu_{it},
\end{equation}
where $\Delta$ denotes the change over a $k$-year horizon, $z_{it}$ is taxable income (wages
minus the applicable standard deduction), $\tau_{it}$ is the combined federal-plus-state
marginal income tax rate from TAXSIM35, $\log z_{i,t-k}$ is the lagged log income base
control (to address mean reversion), and $\alpha_t$ is a period fixed effect. The hat
denotes the fitted value from the first stage.

The first stage regresses the actual change in the log net-of-tax rate on the simulated
change:
\begin{equation}
\label{eq:eti_fs}
\Delta \log(1 - \tau_{it}) = \pi \cdot \Delta \log(1 - \tilde{\tau}_{it}) + \beta' \cdot \log z_{i,t-k} + \alpha'_t + u_{it},
\end{equation}
where $\tilde{\tau}_{it}$ is the marginal tax rate that individual $i$ \emph{would face}
in period $t$ if her income remained at its base-period level inflated by average economy-wide
earnings growth, and only the tax schedule changed. The simulated instrument $\Delta \log(1 -
\tilde{\tau}_{it})$ is correlated with actual rate changes driven by legislated reforms but
uncorrelated with individual income adjustments. The identifying assumption is the exclusion
restriction: $\Delta \log(1 - \tilde{\tau}_{it}) \perp \nu_{it}$, conditional on the income
base control and period effects.

I estimate equation~(\ref{eq:eti}) separately in two periods:
\begin{itemize}
\item \textbf{Annual (1978--1993)}: $k = 3$-year horizon; NLSY79 workers are ages 17--35;
  identifying reforms are ERTA (1981) and TRA (1986).
\item \textbf{Biennial (1995--2019)}: $k = 2$-period horizon; NLSY79 workers are ages 31--62;
  identifying reforms are EGTRRA (2001), JGTRRA (2003), and TCJA (2017).
\end{itemize}

The two-period split is motivated by Corollary~\ref{prop:gamma_arc}: young workers in the
signaling phase optimize reputation accumulation rather than tax liability, so the behavioral
ETI should be near zero in the annual period and positive in the biennial period when workers
have established reputations. This is a testable prediction of the model.

\paragraph{Threats to identification.} Three concerns arise in ETI estimation
\citep{kopczuk2005tax, saez2012elasticity}.

\emph{Mean reversion.} High earners in the base period mechanically show income convergence
in subsequent periods, creating a spurious negative correlation between the level of income
(and hence marginal tax rate) and income growth. The base-period income control $\log z_{i,t-k}$
in equation~(\ref{eq:eti}) addresses this directly, as in \citet{gruber2002elasticity}. A
remaining concern is that the simulated rate change $\Delta \log(1 - \tilde{\tau}_{it})$ is
mechanically larger for high-income workers (who face steeper brackets), and if these workers
have systematically different income growth trends, the exclusion restriction is violated.
The base-period control mitigates this, and the strong first-stage $F$-statistics (2,617
annual, 1,146 biennial) confirm the instrument is not weak.

\emph{Income floor endogeneity.} Workers near the income floor may enter or exit the sample
as tax rates change, contaminating the ITI with an extensive-margin effect. I apply income
floors (\$5,000 annual, \$10,000 biennial) and verify robustness to a lower \$1,000 floor;
the near-zero annual ETI is stable across both thresholds (Section~\ref{sec:results}).

\emph{Marital transitions.} Changes in filing status alter marginal rates independently of
the tax schedule, confounding the instrument. I include indicator variables for marital
status transitions (single-to-married and married-to-single) as controls in both stages.
To verify that the results are not driven by marital changers, I replicate the baseline
estimates excluding all person-years with a marital transition; the ETI estimates are
substantively unchanged.

\subsection{Estimating $\alpha$: Pareto Tail Regression}

The income distribution above the 80th percentile is well-approximated by a Pareto distribution
with tail parameter $\alpha$. I estimate $\alpha$ via log-rank regression: ordering workers by
annual income and regressing $\log(\text{rank})$ on $\log(\text{income})$ for the top 20
percent of the annual distribution:
\begin{equation}
\label{eq:pareto}
\log(\text{rank}_i) = c - \alpha \cdot \log(w_i) + \eta_i.
\end{equation}
The slope coefficient $-\hat{\alpha}$ is the Pareto exponent, and the goodness-of-fit ($R^2$)
measures how well the Pareto distribution approximates the income tail. Standard errors are
computed by OLS with heteroskedasticity-robust standard errors.

%==============================================================================
\section{Results}
\label{sec:results}
%==============================================================================

\subsection{The $\gamma$ Career Arc}

Table~\ref{tab:gamma} reports the career-stage $\gamma$ estimates from equation~(\ref{eq:fe_main}),
separately for five experience quintiles and for the full career. The full-career estimate is
$\hat{\gamma} = 0.937$ (SE $= 0.011$, $t = 89.2$, $N = 179{,}209$): each 10 percent increase
in cumulative hours is associated with a 9.4 percent wage increase within person, conditional
on potential experience and year effects. This elasticity is below 1.0, implying that $\delta <
1 - \varepsilon^w$, consistent with the model's support constraint.

The quintile estimates reveal a career arc. Early-career workers (Q1, pot\_exp 0--8 years)
show $\hat{\gamma} = 1.10$ (SE $= 0.021$); mid-career workers (Q4, pot\_exp 15--25 years)
show $\hat{\gamma} = 0.927$ (SE $= 0.057$). A joint Wald test rejects equality across stages
at $F(4, 12{,}686) = 16.87$ ($p < 0.001$). The Q5 estimate (pot\_exp $> 25$ years) is not
used in the structural analysis: Section~\ref{sec:robustness} documents that Q5 is
non-identified due to survivor selection and a within-$R^2$ of 2.6 percent.

A notable feature of the Q1 estimate is that $\hat{\gamma}_{\text{Q1}} = 1.10 > 1$. Under
the structural model with $\varepsilon^w \approx 0$ for young workers, equation~(\ref{eq:gamma_structural})
implies $\gamma = \delta / (1 - \delta)$, which exceeds 1 only if $\delta > 0.5$. Two
interpretations are consistent with this. First, the polynomial specification---which does not
discretize careers into bins---estimates $\hat{\beta}_0 = 0.974$ at career entry, below 1.0
and within the model's support. The quintile step function attributes an average of the
declining arc to Q1, which includes exp = 0--8 years; the polynomial is more precise about
the entry value. Second, the annual sample for young workers (ages 17--35) spans a period
when the panel is most unbalanced due to school-to-work transitions, and short-panel dynamics
may mildly inflate the Q1 FE estimate. The structural inference relies on the full-career
$\hat{\gamma} = 0.937$ and the polynomial $\hat{\beta}_0 = 0.974$; the Q1 quintile estimate
is reported for completeness but is not used to compute $\hat{\delta}$ or $\hat{\tau}_p$.

The polynomial specification (\ref{eq:fe_poly}) provides a sharper picture of the arc.
Estimates from this specification are:
\begin{align*}
\hat{\beta}_0 &= 0.974 \;\; (SE = 0.011,\; t = 90.9) \\
\hat{\beta}_1 &= -0.0134 \;\; (SE = 0.0015,\; t = -9.26) \\
\hat{\beta}_2 &= +0.0003 \;\; (SE = 0.00004,\; t = 7.07).
\end{align*}
The dominant term is $\hat{\beta}_1 < 0$: $\gamma$ declines approximately linearly at
$-0.013$ per year of experience. The small positive $\hat{\beta}_2$ implies a slight upturn
after the career minimum, but $|\hat{\beta}_2| < |\hat{\beta}_1|/5$, so the quadratic term
is negligible relative to the linear decline. The career minimum is at $e^* = -(-0.0134) /
(2 \times 0.0003) = 21.6$ years, where $\hat{\gamma}(21.6) = 0.829$. Figure~\ref{fig:gamma}
plots the full $\hat{\gamma}(\text{exp})$ arc.

The joint test $H_0\colon \beta_1 = \beta_2 = 0$ is rejected at $F(2, 11524) = 43.20$
($p < 0.001$), rejecting the constant-$\gamma$ Pareto case. This is consistent with
Proposition~\ref{prop:gamma_arc}: employer uncertainty is highest at career entry and declines
as the worker accumulates a public work history.

\paragraph{Selection bias.} Table~\ref{tab:gamma} also reports OLS estimates for comparison.
The full-career OLS elasticity is $\hat{\gamma}_{\text{OLS}} = 1.022$, substantially above
the FE estimate of 0.937, indicating that high-ability workers supply more hours
(positive selection on ability). Within the career arc, the OLS--FE gap reverses sign: in
early career (Q1--Q3), $\hat{\gamma}_{\text{FE}} > \hat{\gamma}_{\text{OLS}}$, while in late
career (Q4) $\hat{\gamma}_{\text{OLS}} > \hat{\gamma}_{\text{FE}}$. This career-arc reversal
in selection bias is consistent with the employer learning mechanism: early-career
cross-sectional wages are compressed because employers pool workers of different abilities,
so the cross-sectional $\gamma$ understates the within-person return. As employers learn
and wages begin to spread with ability, cross-sectional wages eventually become more
dispersed, and the OLS--FE ordering restores. This pattern is not predicted by standard
human capital models and provides additional evidence for the signaling mechanism.

\subsection{The Elasticity of Taxable Income}

Table~\ref{tab:eti} reports the ETI estimates from equation~(\ref{eq:eti}) for both the
annual and biennial periods. The first-stage $F$-statistics are 2,617 (annual) and 1,146
(biennial), both far above the conventional weak-instrument threshold of 10; the instrument
is strong in both periods.

\paragraph{Annual period (1978--1993).} The annual ETI is $\hat{\varepsilon}^w = -0.048$
(SE $= 0.080$, $t = -0.60$, $N = 53{,}632$): statistically indistinguishable from zero.
This estimate is stable under alternative income floors: lowering the threshold from \$5,000
to \$1,000 adds 5,796 near-workers and yields $\hat{\varepsilon}^w = -0.049$ (SE $= 0.082$,
$t = -0.60$, $N = 59{,}428$, $F = 2{,}500$). The near-zero annual ETI is not explained by
weak identification: the first-stage $F$ at the \$1,000 floor is 2,500, confirming that
near-workers face genuine reform-driven variation in their net-of-tax rate, to which they
simply do not respond. This is the pattern predicted by the signaling model: young workers
(ages 17--35) face strong reputation-building incentives that dominate tax optimization.

\paragraph{Biennial period (1995--2019).} The biennial ETI is $\hat{\varepsilon}^w = +0.297$
(SE $= 0.121$, $t = 2.46$, $N = 47{,}201$): statistically significant at the 5 percent level.
This estimate covers prime-career and late-career workers (ages 31--62) who have established
labor market reputations and face the intensive-margin labor supply decision without the
competing signaling incentive. The positive and significant biennial ETI is consistent with
a conventional tax elasticity response emerging after the signaling phase concludes.

\paragraph{Lifecycle interpretation.} The contrast between the two periods—$\hat{\varepsilon}^w
\approx 0$ for young workers, $\hat{\varepsilon}^w \approx 0.30$ for prime-career workers—is
itself a prediction of Corollary~\ref{prop:gamma_arc}. In the signaling phase, the private
return to additional work hours includes both the current wage and the future wage premium
from improved employer beliefs. This signaling component acts as an implicit subsidy to labor
effort that crowds out the tax response; workers supply hours at the margin of reputation
accumulation rather than net-of-tax wages. The biennial ETI of 0.297 is the behavioral
parameter that applies to workers who have exited the signaling phase, and it is this estimate
that enters the $\tau_p$ formula as the structural $\varepsilon^w$.

The biennial ETI of 0.297 is below \citepos{gruber2002elasticity} benchmark estimate of
approximately 0.4 for higher-income taxpayers. This comparison is directionally consistent:
the NLSY79 sample covers a broader income distribution (including workers well below the top
quintile), and lower-income workers are generally found to have smaller ETIs because they
face lower marginal rates and have fewer tax-avoidance margins \citep{saez2012elasticity}.
The sign (positive) and rough order of magnitude (below 0.5) are consistent with the prior
literature on prime-career workers.

\subsection{The Optimal Pigouvian Subsidy}

\paragraph{Pareto tail.} The log-rank regression yields $\hat{\alpha} = 3.295$
(SE $= 0.014$, $R^2 = 0.985$). The high $R^2$ confirms that the Pareto distribution
approximates the top 20 percent of the NLSY79 income distribution well. The estimate is
below the commonly referenced value of $\alpha \approx 1.5$--$2.0$ for the full US income
distribution from tax records \citep{saez2001using}, but recall that $\alpha$ estimates
depend heavily on the income concept and population. Earnings-only Pareto estimates for
mid-career cohorts tend to be substantially higher (less heavy-tailed) than estimates from
household income or wealth distributions, because the earnings distribution among prime-age
workers is compressed relative to the total income distribution which includes capital income
and retirement withdrawals. \citet{diamond1998optimal} survey Pareto exponents in the range
$\alpha \in [1.5, 4]$ for various income concepts and time periods; our $\hat{\alpha} = 3.295$
falls in the upper part of this range, consistent with an earnings-only, prime-age cohort.

\paragraph{Structural $\delta$.} Applying equation~(\ref{eq:delta}) with the full-career
FE estimate $\hat{\gamma} = 0.937$ and the biennial ETI $\hat{\varepsilon}^w = 0.297$:
\begin{equation}
\hat{\delta} = \frac{0.937 \times 1.297}{1.937} = 0.628.
\end{equation}
This signaling parameter implies that after a given period, there remains a 62.8 percent
probability that employers have not yet fully resolved uncertainty about a worker's type.
For comparison, \citepos{sztutman2024} HRS estimate implies $\delta \approx 0.17$ for
workers ages 50 and above; the NLSY79 cohort, which spans ages 17--62, has substantially
higher residual employer uncertainty on average.

\paragraph{Pigouvian tax.} Applying equation~(\ref{eq:taup}):
\begin{equation}
\hat{\tau}_p = \frac{\hat{\delta}}{\hat{\alpha}} = \frac{0.628}{3.295} = 19.1\%.
\end{equation}
This is the central estimate of the optimal Pigouvian labor income subsidy. Table~\ref{tab:taup}
summarizes the components and several sensitivity variants.

\paragraph{Confidence interval.} The dominant source of uncertainty is the biennial ETI
$\hat{\varepsilon}^w$ (SE $= 0.121$). Propagating uncertainty through
$\delta = \gamma(1+\varepsilon^w)/(1+\gamma)$ and $\tau_p = \delta / \alpha$ via the
delta method:
\begin{equation*}
\text{SE}(\hat{\tau}_p) = \frac{1}{\hat{\alpha}} \sqrt{
\left(\frac{\partial \hat{\delta}}{\partial \varepsilon^w}\right)^2 \text{SE}(\hat{\varepsilon}^w)^2 +
\left(\frac{\partial \hat{\delta}}{\partial \gamma}\right)^2 \text{SE}(\hat{\gamma})^2
} \approx 0.0178,
\end{equation*}
where $\partial \hat{\delta}/\partial \varepsilon^w = \hat{\gamma}/(1+\hat{\gamma}) = 0.484$
and $\partial \hat{\delta}/\partial \gamma = (1+\hat{\varepsilon}^w)/(1+\hat{\gamma})^2 = 0.346$.
The $\hat{\varepsilon}^w$ term contributes 99.6 percent of the total variance; the $\hat{\gamma}$
term contributes 0.4 percent.

The 95 percent confidence interval is $\hat{\tau}_p \pm 1.96 \times 0.0178 = [15.6\%, 22.6\%]$.
A clustered bootstrap (B $= 500$, resampling persons) that re-estimates $\hat{\varepsilon}^w$
on each resample and propagates through the formula yields an identical interval of
$[15.6\%, 22.5\%]$, validating the delta-method approximation. A joint bootstrap (B $= 200$)
that simultaneously re-estimates both $\hat{\gamma}$ and $\hat{\varepsilon}^w$ on each resample
yields a slightly narrower interval of $[15.7\%, 20.8\%]$, confirming that uncertainty in
$\hat{\gamma}$ does not widen the interval—consistent with the analytic result that $\hat{\gamma}$
contributes less than 1 percent of the variance in $\hat{\tau}_p$.

The primary reported interval is the delta-method CI of $[15.6\%, 22.6\%]$, which is the
widest and most conservative of the three approaches.

%==============================================================================
\section{Cross-Validation of the Employer Learning Mechanism}
\label{sec:validation}
%==============================================================================

The structural estimate $\hat{\tau}_p = 19.1\%$ rests on the interpretation of $\hat{\gamma}$
as an employer-learning parameter. Two independent tests of this interpretation are available
in the NLSY79: sector heterogeneity in $\gamma$ and the Altonji--Pierret (AP) test for
gradual ability revelation. Both tests use only the NLSY79 sample period 1979--1993 when
occupation codes are available.

\subsection{Sector Heterogeneity in $\gamma$}
\label{sec:hetero}

If $\gamma$ reflects employer learning, it should be higher in sectors where output quality
is less directly observable and employer learning is more important. I assign NLSY79 occupations
(1-digit) to three signaling-intensity groups: Low (farm workers, laborers: occupations 7--8),
Medium (clerical, sales, operatives, craftsmen, service: occupations 3--6, 9), and High
(professional/technical, managers: occupations 1--2). The predicted ordering is
$\gamma_{\text{Low}} < \gamma_{\text{Med}} < \gamma_{\text{High}}$.

Table~\ref{tab:sector} reports FE estimates of $\gamma$ by group. The estimates are
$\hat{\gamma}_{\text{Low}} = 0.900$ (SE $= 0.036$), $\hat{\gamma}_{\text{Med}} = 0.953$
(SE $= 0.028$), and $\hat{\gamma}_{\text{High}} = 1.037$ (SE $= 0.027$), a monotone increase
consistent with the signaling prediction. The gap between High and Low groups is 0.137 (95\%
CI: $[0.067, 0.207]$), statistically significant and economically meaningful.
Appendix Table~\ref{tab:gamma_occ} reports $\hat{\gamma}$ for each of the nine 1-digit
occupations underlying these group averages. The range from $\hat{\gamma} = 0.470$ for Farm
Workers to $\hat{\gamma} = 1.235$ for Sales workers---a spread of 0.765---substantially
exceeds the group-level High--Low gap of 0.137, confirming that occupational heterogeneity
in employer learning intensity is pervasive and not merely an artifact of the three-group
aggregation.

\subsection{The Altonji--Pierret Test for Employer Learning}
\label{sec:ap_test}

\citet{altonji2001employer} show that under employer learning, the wage premium for observable
ability proxies (such as test scores) should \emph{increase} with experience, as market-wide
wages become more responsive to ability as employers learn. I implement their test by
interacting AFQT percentile scores (measured pre-labor-market in 1980) with potential
experience in a wage regression:
\begin{equation}
\label{eq:ap}
\log w_{it} = \alpha + \phi_1 \cdot \text{AFQT}_i + \phi_2 \cdot (\text{AFQT}_i \times \text{exp}_{it})
+ \beta \cdot X_{it} + \varepsilon_{it}.
\end{equation}
The employer-learning prediction is $\phi_2 > 0$: the AFQT premium rises with experience.

Averaging across all occupations, the estimated AP slope is $\hat{\phi}_2 = 0.00960$ per year
for the medium-signaling group ($t = 4.62$), consistent with gradual ability revelation. The
full-sample group ordering is non-monotone (Low $= 0.00511$, Medium $= 0.00960$, High $= 0.00666$),
driven by a composition artifact: Laborers in the Low group have a spuriously high AP slope
of 0.01316 when pooled across tenure levels, because short-tenure workers dominate the sample
and their early-career AFQT gradient conflates ability revelation with initial wage-setting.

Restricting to workers with at least 5 years of occupation tenure, the group-level AP slopes are
Low $= 0.00432$ and High $= 0.00882$, reversing the ordering and delivering the signaling
prediction. The pattern is anchored by two occupations: Managers/Admin (the established
High-signaling occupation) shows 0.01493 ($t = 2.85$, $N = 3{,}409$), while Laborers collapse
from 0.01316 to 0.00015 ($t = 0.01$, $N = 2{,}292$)---a near-zero estimate consistent with
the absence of employer learning where output is directly observable. The Low $<$ High ordering
in Table~\ref{tab:sector} is consistent with the signaling mechanism but should be interpreted
with caution: the 1979--1993 window restricts all workers to early career stages, and
Professional/Technical workers within the High group remain insignificant ($\hat{\phi}_2 =
0.00271$, $t = 0.62$), possibly reflecting genuine occupational heterogeneity in learning speed
within the manager/professional category.

\subsection{Horse Race: Cumulative versus Recent Hours}

If cumulative hours signal ability through their public visibility to all employers, then
the wage return to cumulative hours should exceed the return to recent hours early in the
career (when the cumulative record is informative) but not late in the career (when
reputation is already established and recent effort is the marginal signal). I test this
by jointly estimating:
\begin{equation}
\label{eq:horse}
\log w_{it} = \alpha_i + \gamma_c \cdot \log H_{it} + \gamma_r \cdot \log h_t
+ \text{controls} + \varepsilon_{it},
\end{equation}
where $h_t$ is hours worked in the current year. The signaling prediction is
$\gamma_r > \gamma_c$ early in the career, converging to $\gamma_c > \gamma_r$ late in the
career as reputations are established.

In the early career (Q1), $\hat{\gamma}_r = 0.815$ and $\hat{\gamma}_c = 0.166$---recent
hours dominate by a ratio of 4.9 to 1, consistent with the signaling interpretation that
current-period face-time is the principal signal. By mid-career (Q4), $\hat{\gamma}_r = 0.529$
and $\hat{\gamma}_c = 0.396$---the two are roughly comparable. This transition from recent-hours
dominance to near-parity is consistent with the career arc of Proposition~\ref{prop:gamma_arc}
and further corroborates the employer learning mechanism.

%==============================================================================
\section{Robustness}
\label{sec:robustness}
%==============================================================================

\subsection{Collinearity: Frisch--Waugh--Lovell Verification}

A concern in equation~(\ref{eq:fe_main}) is that $\log H_{it}$ and $\text{pot\_exp}_{it}$
are mechanically correlated within person: both grow over the career. I quantify the
collinearity as follows: regressing $\log H_{it}$ on $\text{pot\_exp}_{it}$, $\text{pot\_exp}_{it}^2$,
and year dummies within each person and computing the $R^2$ yields 71.6 percent---so 28.4
percent of within-person cumulative-hours variation is orthogonal to the experience controls.
This residual variation identifies $\gamma$.

The FWL theorem guarantees that the FE estimator applied to the full regression produces the
same $\hat{\gamma}$ as a two-stage procedure that first residualizes $\log H_{it}$ and $\log
w_{it}$ on the controls, then regresses the wage residual on the hours residual. I verify this
numerically: the difference between $\hat{\gamma}$ from the full regression and from the FWL
procedure is exactly 0.0000 for every career stage. This confirms that the estimated career
arc---including the Q1 elevation and Q4 trough---is not an artifact of collinearity.
Collinearity inflates standard errors (reducing precision) but does not bias point estimates.

\subsection{Simultaneity: Instrumental Variables for $\gamma$}

Within-year simultaneity is a second concern: a high job-match year may generate both more
hours and higher wages simultaneously, creating upward bias in $\hat{\gamma}$. To address
this, I instrument $\log H_{it}$ with its one-year lag, $\log H_{i,t-1}$. Lagged cumulative
hours are highly correlated with current cumulative hours ($\hat{\beta}_{\text{FS}} = 0.670$,
first-stage $t = 169$, within-$R^2 = 0.980$) but uncorrelated with the current-period
match-quality shock, conditional on worker fixed effects and year effects.

The IV estimate is $\hat{\gamma}_{\text{IV}} = 0.697$ (SE $= 0.007$, $t = 103$), compared
to the FE estimate of 0.937. The gap of 0.239 estimates the within-year simultaneity bias.
This moderate-to-large bias implies that the FE estimate is an upper bound on the true $\gamma$;
the IV estimate of 0.697 provides a lower bound. Both estimates imply $\delta < 1$
(consistent with the model), and both imply $\tau_p > 15\%$: applying $\hat{\gamma}_{\text{IV}} = 0.697$
with $\hat{\varepsilon}^w = 0.297$ yields $\hat{\delta} = 0.533$ and $\hat{\tau}_p = 0.533/3.295 = 16.2\%$.
The structural parameter range implied by FE and IV is $\hat{\tau}_p \in [16.2\%, 19.1\%]$.

\subsection{Heterogeneous Trends: First-Difference Estimator}

Person fixed effects absorb time-invariant heterogeneity but not person-specific linear
trends in wages. If some workers have both faster wage growth and faster hours growth (e.g.,
due to unobserved career investment), the FE estimate of $\gamma$ could be upward biased.
I address this by first-differencing the wage equation:
\begin{equation}
\Delta \log w_{it} = \gamma \cdot \Delta \log H_{it} + \beta \cdot \Delta \text{exp}_{it}
+ \Delta \lambda_t + \Delta \varepsilon_{it},
\end{equation}
which removes both fixed effects and person-specific linear trends. The first-difference
estimate is $\hat{\gamma}_{\text{FD}} = 1.367$ (SE $= 0.027$)---substantially \emph{above}
the FE estimate of 0.937, not below. The bias is negative throughout all career stages
(FD $>$ FE), ruling out the heterogeneous upward trend explanation.

The FD estimate is above the FE estimate because annual first-differencing identifies the
return to the \emph{marginal} resume signal: $\Delta \log H_{it} \approx h_t / H_t$ (this
year's hours as a fraction of career total), which captures within-year employment-intensity
shocks that covary with match-quality and wages. The FE estimator identifies $\gamma$ from
the full within-person career trajectory, averaging over year-to-year shocks, and is the
appropriate structural parameter for the career-arc model. The key diagnostic from the FD
comparison is negative: the Q4 trough in $\gamma$ is confirmed in both estimators (FD-Q4 $=
1.046$ versus FD-Q1 $= 1.398$), ruling out heterogeneous trends as the explanation for the
career-arc decline.

\subsection{Q5 Non-Identification}

The Q5 quintile estimate (pot\_exp $> 25$ years) is $\hat{\gamma}_{\text{Q5}} = 1.115$
(SE $= 0.110$, 95\% CI: $[0.900, 1.330]$). This CI contains every other quintile estimate
and the Q5 estimate is statistically indistinguishable from all of Q1--Q4. Two compounding
problems prevent Q5 identification. First, the within-$R^2$ in the Q5 subsample is 2.6
percent---employer learning essentially explains none of the within-person wage variation
in late career, consistent with employer uncertainty having largely resolved. Second, workers
who survive into Q5 (pot\_exp $> 25$ years) have mean Q4 log wages 0.567 points above workers
who exit at Q4---a massive positive survivor selection that mechanically inflates the Q5
estimate. For these reasons, Q5 is excluded from the structural career arc. The reported
career-stage estimates cover Q1--Q4 (pot\_exp 0--25 years, ages roughly 22--47).

\subsection{Sensitivity to $\alpha$}

The income distribution's Pareto tail is assumed to hold at $\alpha = 3.295$ as estimated.
Pareto tail estimates are sensitive to the threshold above which the Pareto distribution is
assumed to hold and to the specific dataset. Table~\ref{tab:taup} reports $\hat{\tau}_p$
under three alternative values: $\alpha = 1.5$ (very heavy tail), $\alpha = 2.5$ (moderate),
and $\alpha = 3.0$ (close to our estimate). At $\alpha = 1.5$, $\hat{\tau}_p = 41.9\%$; at
$\alpha = 2.5$, $\hat{\tau}_p = 25.1\%$; at $\alpha = 3.0$, $\hat{\tau}_p = 20.9\%$. The
sensitivity to $\alpha$ is the primary source of model uncertainty outside the estimated
CI. Our direct $\hat{\alpha} = 3.295$ (R$^2 = 0.985$) is close to the $\alpha = 3.0$
benchmark, and the robustness range $[20.9\%, 41.9\%]$ for $\alpha \in [1.5, 3.0]$ shows
that the main message---$\tau_p$ is economically significant---holds throughout.

%==============================================================================
\section{Discussion}
\label{sec:discussion}
%==============================================================================

\subsection{The Lifecycle Arc of $\delta$}

The structural signaling parameter $\hat{\delta} = 0.628$ for the NLSY79 cohort (average
age 33) is roughly 3.7 times \citepos{sztutman2024} HRS estimate of $\delta \approx 0.17$
for workers ages 50 and above. This factor-of-four discrepancy is predicted by the model.

The $\delta$ lifecycle arc can be traced using the polynomial $\hat{\gamma}(\text{exp})$ arc
estimated in Section~\ref{sec:results}. Applying equation~(\ref{eq:delta}) at each experience
level with fixed $\hat{\varepsilon}^w = 0.297$:
\[
\hat{\delta}(\text{exp}) = \frac{\hat{\gamma}(\text{exp}) \times 1.297}{1 + \hat{\gamma}(\text{exp})}.
\]
This yields $\hat{\delta}(0) = 0.640$ at career entry, declining to $\hat{\delta}(20) = 0.588$
at the midpoint, and $\hat{\delta}(40) = 0.627$ at the end of the NLSY79 sample window
(where the slight upturn reflects the survivor selection discussed in Section~\ref{sec:robustness}).
Table~\ref{tab:delta_arc} reports these values alongside the implied $\hat{\tau}_p(\text{exp})$.

A notable finding from Table~\ref{tab:delta_arc} is that $\hat{\tau}_p(\text{exp})$ is
remarkably stable within each ETI regime: it varies by only 1.6 percentage points across
experience levels under the biennial ETI (from 17.8\% to 19.4\%) and by 1.3 percentage
points under the annual ETI. The primary source of variation in the optimal subsidy is not
the career stage within the NLSY79 sample but the transition between lifecycle ETI regimes:
the biennial estimate (19.1\%) exceeds the annual estimate (14.0\%) by 5.1 percentage points,
reflecting the regime change in behavioral response as careers mature.

The comparison to \citepos{sztutman2024} HRS estimate ($\tau_p \approx 5\%$) implies a
dramatic reduction in the signaling externality between the career stages covered by the two
datasets. Under the model, this reduction reflects the resolution of employer uncertainty:
by ages 50--70, the labor market has had 30--40 years to learn about worker types, and the
residual signaling externality is correspondingly small. The 73 percent decline in $\delta$
from $0.628$ (NLSY79 average) to $0.17$ (HRS) is consistent with the employer learning
dynamics that Corollary~\ref{prop:gamma_arc} implies.

\paragraph{The appropriate $\varepsilon^w$ for the structural formula.}
Table~\ref{tab:delta_arc} presents two scenarios—using the annual ETI ($\hat{\varepsilon}^w
\approx 0$) and the biennial ETI ($\hat{\varepsilon}^w = 0.297$)—rather than a single
lifecycle estimate. The choice matters because the structural formula $\tau_p = \delta/\alpha$
applies when workers face a genuine tax margin: the Pigouvian subsidy corrects the wedge
between the private and social return to labor, but only insofar as workers respond to the
net-of-tax wage at all. For young workers in the signaling phase (annual period, ages 17--35),
the near-zero annual ETI reflects the dominance of reputation-building incentives over tax
optimization, not a preference for corner solutions. In this regime, the Pigouvian case for
a subsidy is strongest (high $\delta$) but the tax margin is inoperative. The biennial ETI
of 0.297 applies to workers who have exited the signaling phase and face the standard
intensive-margin labor supply decision; it is this estimate that yields the primary $\hat{\tau}_p
= 19.1\%$. A lifecycle-appropriate combination—using annual $\hat{\varepsilon}^w$ for
$\text{exp} \lesssim 15$ and biennial $\hat{\varepsilon}^w$ for $\text{exp} \gtrsim 15$—yields
$\hat{\tau}_p \approx 13$--$14\%$ in early career and $\hat{\tau}_p \approx 18$--$19\%$ for
prime-career workers, with the step at the ETI regime transition (approximately age 30)
as the primary source of variation. The career-length average is approximately 17 percent of
labor income.

\paragraph{Reconciling the flat $\delta$ arc with the employer learning literature.}
The estimated $\hat{\delta}(\text{exp})$ arc declines by only 7 percent from $0.640$ at
career entry to $0.588$ at the mid-career minimum (Table~\ref{tab:delta_arc}), which may
appear inconsistent with \citepos{lange2007speed} finding that employer learning is
substantially complete within 3 years of labor market entry. Two observations reconcile
these findings. First, Lange measures learning from the rising cross-sectional wage premium
to AFQT scores—the speed at which the market price of ability converges from a pooling wage
to its fully revealed level. This object is distinct from the structural $\delta$ in
equation~(\ref{eq:gamma_structural}), which measures the fraction of the wage-hours elasticity
that flows through the signaling channel. Even after the market assigns most of the ability
premium correctly, workers may continue to supply hours partly for reputation reasons if
there remains residual uncertainty about exact productivity type; $\delta$ captures this
sustained informational margin rather than the initial speed of belief updating. Second,
Lange's 3-year result applies precisely to the early-career window where our polynomial
$\hat{\gamma}(\text{exp})$ exhibits its steepest decline---from 0.974 at entry to 0.915 by
$\text{exp} = 5$ and 0.871 by $\text{exp} = 10$. The implied $\hat{\delta}$ declines from
0.640 to 0.604 over the same 10 years (a 6 percent drop), concentrated exactly where Lange
predicts. The flat profile at mid-career (exp 10--25) reflects slow residual learning, not
a contradiction of the rapid initial learning that Lange documents.

\subsection{Policy Implications}

These estimates suggest that the optimal marginal labor income tax rate for prime-career
workers is approximately 19 percentage points below the rate implied by equity and revenue
considerations alone. This finding is consistent with the theoretical arguments of
\citet{boskin1975taxation} and \citet{sztutman2024} that signaling externalities warrant
a Pigouvian labor income subsidy, and it quantifies the magnitude of that subsidy for a
cohort-representative US panel.

Several caveats apply. First, the NLSY79 cohort was born 1957--1964 and may not be
representative of workers in other cohorts or institutional contexts; the lifecycle
dynamics of $\delta$ could differ for cohorts who entered the labor market during periods
of greater wage compression or different sectoral composition. Second, the structural
formula $\tau_p = \delta / \alpha$ assumes the Pareto case and a specific information
structure; departures from these assumptions would change the formula, though Proposition~1
shows that the sign of the externality---and hence the direction of the subsidy---is
robust to the distributional assumption. Third, the estimates rely on the biennial ETI
of $+0.297$; if this estimate reflects mean reversion or other biases not fully addressed
by the Gruber--Saez controls, the implied $\tau_p$ would be smaller. Fourth, the analysis
abstracts from heterogeneity in $\tau_p$ across income groups, which the model predicts
should exist (higher-income workers may have larger signaling externalities). Our estimate
is a population average.

These findings are consistent with, rather than establishing, the case for a Pigouvian
labor income subsidy. A full welfare analysis would require modeling the equity costs of
the subsidy and the distributional incidence of the tax reforms used for identification,
which is beyond the scope of this paper.

%==============================================================================
\section{Conclusion}
\label{sec:conclusion}
%==============================================================================

The wages of NLSY79 workers grow with cumulative hours at an elasticity that declines from
0.97 at labor market entry to 0.83 by midcareer---a pattern that indicates employer uncertainty
about worker ability resolves progressively over careers. This employer learning process
generates a positive externality to labor effort: each additional hour of work creates
information about ability that benefits future employers who did not pay for the additional
labor. The Sztutman (2024) model formalizes this externality and shows that the optimal
Pigouvian labor income subsidy equals $\tau_p = \delta / \alpha$, where $\delta$ is the
signaling probability and $\alpha$ is the Pareto income tail parameter.

I estimate all three sufficient statistics---$\gamma$, $\varepsilon^w$, and $\alpha$---from
the NLSY79 panel and find $\hat{\tau}_p = 19.1\%$ (95\% CI: $[15.6\%, 22.6\%]$). This
estimate is substantially larger than prior estimates for older workers \citep{sztutman2024},
consistent with the prediction that the signaling externality declines over careers as
employer uncertainty resolves. Cross-validation through sector heterogeneity in $\gamma$,
the Altonji--Pierret ability-revelation test, and the horse race between cumulative and
recent hours all support the employer learning interpretation.

Three results stand out. First, the career arc of $\gamma$---declining from 0.97 to 0.83
at career midpoint and confirmed by both FE and FD estimators---provides direct evidence
for the dynamic signaling mechanism in a representative US panel. Second, the lifecycle
pattern in the ETI---near zero for young workers, positive for prime-career workers---is
consistent with a signaling-to-optimization transition over careers and suggests that
standard ETI estimates that pool all career stages may be measuring a mixture of behavioral
responses. Third, the implied signaling parameter $\hat{\delta} = 0.628$ for prime-career
US workers implies that, after a given year, there remains roughly a two-in-three chance
that employers have not fully resolved uncertainty about a worker's type---a substantial
residual uncertainty that justifies a meaningful Pigouvian correction.

These findings raise an open question for future research: how does the optimal Pigouvian
subsidy vary across income groups and occupations? The sector heterogeneity results suggest
that high-signaling occupations have substantially higher $\gamma$ values, which would imply
higher $\delta$ and larger Pigouvian corrections for those workers. A full distributional
analysis of the optimal labor income tax structure, accounting for both the equity objectives
and the occupation-specific signaling externalities, is a promising direction.

%==============================================================================
% REFERENCES
%==============================================================================
\newpage
\bibliographystyle{aer}
\begin{thebibliography}{99}

\bibitem[Altonji and Pierret(2001)]{altonji2001employer}
Altonji, J.~G., and C.~R. Pierret. 2001.
\newblock Employer learning and statistical discrimination.
\newblock \emph{Quarterly Journal of Economics} 116 (1): 313--350.

\bibitem[Boskin(1975)]{boskin1975taxation}
Boskin, M.~J. 1975.
\newblock Efficiency aspects of the differential tax treatment of market and
  household economic activity.
\newblock \emph{Journal of Public Economics} 4 (1): 1--25.

\bibitem[Chetty(2009)]{chetty2009sufficient}
Chetty, R. 2009.
\newblock Sufficient statistics for welfare analysis: A bridge between
  structural and reduced-form methods.
\newblock \emph{Annual Review of Economics} 1: 451--488.

\bibitem[Diamond(1998)]{diamond1998optimal}
Diamond, P. 1998.
\newblock Optimal income taxation: An example with a U-shaped pattern of
  optimal marginal tax rates.
\newblock \emph{American Economic Review} 88 (1): 83--95.

\bibitem[Feenberg and Coutts(1993)]{feenberg1993introduction}
Feenberg, D.~R., and E.~Coutts. 1993.
\newblock An introduction to the TAXSIM model.
\newblock \emph{Journal of Policy Analysis and Management} 12 (1): 189--194.

\bibitem[Farber and Gibbons(1996)]{farber1996learning}
Farber, H.~S., and R.~Gibbons. 1996.
\newblock Learning and wage dynamics.
\newblock \emph{Quarterly Journal of Economics} 111 (4): 1007--1047.

\bibitem[Feldstein(1995)]{feldstein1995effect}
Feldstein, M. 1995.
\newblock The effect of marginal tax rates on taxable income: A panel study of
  the 1986 Tax Reform Act.
\newblock \emph{Journal of Political Economy} 103 (3): 551--572.

\bibitem[Gibbons and Waldman(1999)]{gibbons1999careers}
Gibbons, R., and M.~Waldman. 1999.
\newblock A theory of wage and promotion dynamics inside firms.
\newblock \emph{Quarterly Journal of Economics} 114 (4): 1321--1358.

\bibitem[Gruber and Saez(2002)]{gruber2002elasticity}
Gruber, J., and E.~Saez. 2002.
\newblock The elasticity of taxable income: Evidence and implications.
\newblock \emph{Journal of Public Economics} 84 (1): 1--32.

\bibitem[Kopczuk(2005)]{kopczuk2005tax}
Kopczuk, W. 2005.
\newblock Tax bases, tax rates and the elasticity of reported income.
\newblock \emph{Journal of Public Economics} 89 (11-12): 2093--2119.

\bibitem[Lange(2007)]{lange2007speed}
Lange, F. 2007.
\newblock The speed of employer learning.
\newblock \emph{Journal of Labor Economics} 25 (1): 1--35.

\bibitem[Mirrlees(1971)]{mirrlees1971exploration}
Mirrlees, J.~A. 1971.
\newblock An exploration in the theory of optimum income taxation.
\newblock \emph{Review of Economic Studies} 38 (2): 175--208.

\bibitem[Saez(2001)]{saez2001using}
Saez, E. 2001.
\newblock Using elasticities to derive optimal income tax rates.
\newblock \emph{Review of Economic Studies} 68 (1): 205--229.

\bibitem[Saez, Slemrod, and Giertz(2012)]{saez2012elasticity}
Saez, E., J.~Slemrod, and S.~H. Giertz. 2012.
\newblock The elasticity of taxable income with respect to marginal tax rates:
  A critical review.
\newblock \emph{Journal of Economic Literature} 50 (1): 3--50.

\bibitem[Spence(1973)]{spence1973job}
Spence, M. 1973.
\newblock Job market signaling.
\newblock \emph{Quarterly Journal of Economics} 87 (3): 355--374.

\bibitem[Stiglitz(1982)]{stiglitz1982self}
Stiglitz, J.~E. 1982.
\newblock Self-selection and Pareto efficient taxation.
\newblock \emph{Journal of Public Economics} 17 (2): 213--240.

\bibitem[Sztutman(2024)]{sztutman2024}
Sztutman, A. 2024.
\newblock Dynamic job market signaling and optimal taxation.
\newblock Working paper, Carnegie Mellon University Tepper School of Business.

\end{thebibliography}

%==============================================================================
% APPENDIX: TABLES AND FIGURES
%==============================================================================
\newpage
\appendix

%==============================================================================
\section{Proofs}
\label{app:proofs}
%==============================================================================

\subsection*{Proof of Proposition~\ref{prop:taup}}

\noindent \textit{Statement.} Under Assumptions~\ref{ass:pareto}--\ref{ass:learning},
the optimal Pigouvian labor income subsidy is $\tau_p = \delta / \alpha$.

\medskip
\noindent \textit{Proof} (following \citealt{sztutman2024}, Proposition~3; all steps
self-contained).

\medskip
\noindent \textbf{Step 1: Two Pareto properties used in the proof.}

\noindent\textit{(a) Constant inverse hazard rate.} Under Assumption~\ref{ass:pareto},
$1-F(\theta) = (\underline{\theta}/\theta)^\alpha$ and $f(\theta) = \alpha
\underline{\theta}^\alpha \theta^{-(\alpha+1)}$. Computing the inverse hazard:
\begin{equation}
\label{eq:pareto_hazard}
\frac{1-F(\theta)}{\theta\,f(\theta)}
= \frac{(\underline{\theta}/\theta)^\alpha}{\;\theta \cdot \alpha\underline{\theta}^\alpha
  \theta^{-(\alpha+1)}\;}
= \frac{\underline{\theta}^\alpha\,\theta^{-\alpha}}{\alpha\,\underline{\theta}^\alpha\,
  \theta^{-\alpha}}
= \frac{1}{\alpha}.
\end{equation}
This ratio equals $1/\alpha$ at \emph{every} $\theta \geq \underline{\theta}$. This is the
\emph{constant inverse hazard property} of the Pareto distribution, which no other distribution
possesses and which is responsible for the scale-free form of the optimal tax formula.

\medskip
\noindent\textit{(b) Posterior mean under Pareto.} In the signaling equilibrium of
Assumption~\ref{ass:learning}, wages equal the market's posterior mean conditional on
cumulative hours: $w(H) = \mathbb{E}[\theta \mid H_t = H]$. Let $\underline{\theta}(H)$ denote
the minimum type consistent with having accumulated at least $H$ hours in equilibrium---a
threshold that is strictly increasing in $H$ by the monotone-belief property of
Assumption~\ref{ass:learning}. The market conditions on $\theta \geq \underline{\theta}(H)$, so:
\begin{equation}
\label{eq:pareto_posterior}
w(H) = \mathbb{E}[\theta \mid \theta \geq \underline{\theta}(H)]
= \int_{\underline{\theta}(H)}^{\infty} \theta \cdot
  \frac{f(\theta)}{1-F(\underline{\theta}(H))}\,d\theta.
\end{equation}
Substituting $f(\theta) = \alpha\underline{\theta}^\alpha\theta^{-(\alpha+1)}$ and
$1-F(\underline{\theta}(H)) = (\underline{\theta}/\underline{\theta}(H))^\alpha$:
\begin{align*}
w(H)
&= \int_{\underline{\theta}(H)}^{\infty} \theta \cdot
   \frac{\alpha\underline{\theta}^\alpha\,\theta^{-(\alpha+1)}}{(\underline{\theta}/
   \underline{\theta}(H))^\alpha}\,d\theta
= \frac{\alpha\,\underline{\theta}(H)^\alpha}{\alpha - 1}\cdot
  \bigl[\underline{\theta}(H)\bigr]^{-(\alpha-1)} \\
&= \underline{\theta}(H)\cdot\frac{\alpha}{\alpha-1}.
\end{align*}
The wage at any signal level $H$ equals the minimum plausible type scaled by the constant
$\alpha/(\alpha-1)$. Equivalently, $\underline{\theta}(H) = w(H)\cdot(\alpha-1)/\alpha$, and
the ``signaling premium'' above the minimum type---the share of the wage due to favorable
beliefs---equals $1/\alpha$ of the wage at every income level:
\begin{equation}
\label{eq:signaling_premium}
\frac{w(H) - \underline{\theta}(H)}{w(H)} = \frac{1}{\alpha}.
\end{equation}

\medskip
\noindent \textbf{Step 2: Private vs.\ social incentives.}

Each worker chooses hours $h_t$ to maximize $(1-\tau)\,w(H_t) - v(h_t)$, taking the wage
function $w(\cdot)$ as given. The private first-order condition is $(1-\tau)\,\partial w/\partial H = v'(h)$.

Workers fail to internalize the \emph{informational externality}. When worker $i$ increases
$H_i$, this shifts the threshold $\underline{\theta}(H)$ upward at $H = H_i$, raising the
posterior mean for every worker whose observable history coincides with $H_i$. The wage gain
$dw(H_i) = (\alpha/(\alpha-1))\,d\underline{\theta}$ (from equation~\eqref{eq:pareto_posterior})
accrues to worker $i$ privately, but the same upward belief revision benefits all other workers
at signal level $H_i$ as well. The social planner internalizes this externality through a
Pigouvian subsidy.

Of the total wage gain from an additional unit of cumulative hours, fraction $\delta$ is
attributable to employer learning (the signaling channel) and fraction $(1-\delta)$ to genuine
human capital accumulation. Only the $\delta$ component generates an externality, because only
employer-learning wage gains are belief-driven: the social benefit of improved beliefs propagates
to all workers at the same signal level, not just the worker whose hours triggered the update.

\medskip
\noindent \textbf{Step 3: The augmented Pigouvian first-order condition.}

The social planner's optimum satisfies the Saez (2001) perturbation condition augmented by
the informational externality \citep{saez2001using, sztutman2024}. The first-order condition
for the optimal Pigouvian subsidy at income level $w$ is:
\begin{equation}
\label{eq:pigouvian_intermediate}
\tau_p \;=\; \Omega \cdot \frac{1 - F(\theta(w))}{w \cdot f(\theta(w))},
\end{equation}
where $\Omega$ is the \emph{per-unit externality}: the social welfare gain per dollar of labor
income arising from the positive informational spillover of labor supply at income level $w$.
The factor $[1-F(\theta(w))]/[w\,f(\theta(w))]$ is the inverse hazard rate, which in the
Saez framework measures the fiscal cost of correcting the externality through the tax schedule
(a larger right tail means more inframarginal workers, so the same fiscal transfer covers more
people relative to the behavioral distortion at the margin).

To understand the three terms in the social planner's optimality condition, consider a virtual
reform that cuts the marginal tax rate by $d\tau$ on all incomes near $w$:

\emph{(a) Mechanical fiscal cost.} Revenue falls by $d\tau$ per worker above $w$, totaling
$d\tau\cdot[1-F(\theta(w))]\cdot w\,f(w)\,dw$ in present value over the income distribution.

\emph{(b) Behavioral revenue gain.} Workers at $w$ increase hours by
$\varepsilon_c\,d\tau/(1-\tau)$; this generates additional income and partially offsets the
mechanical cost. This behavioral term governs $t_{\text{equity}}$ (the optimal redistributive
rate absent externalities); it does not enter the Pigouvian correction $\tau_p$.

\emph{(c) Externality gain.} The behavioral hours increase at $w$ raises $\underline{\theta}(H(w))$,
shifting the posterior $\mathbb{E}[\theta \mid H(w)]$ upward and raising wages for all workers at
signal level $H(w)$. The social planner values this belief update insofar as it reduces the
under-provision of labor caused by the informational externality. By
Assumption~\ref{ass:learning} and the structural decomposition of
Section~\ref{sec:model}, fraction $\delta$ of any wage gain is attributable to employer
learning; this is the belief-driven component that generates positive externalities for other
workers. Accordingly, $\Omega = \delta$: the per-unit externality equals the employer learning
fraction.

\medskip
\noindent \textbf{Step 4: Applying the Pareto inverse hazard.}

Substituting $\Omega = \delta$ and the constant inverse hazard~\eqref{eq:pareto_hazard}
into~\eqref{eq:pigouvian_intermediate}:
\begin{equation}
\tau_p = \delta \cdot \frac{1-F(\theta(w))}{w\,f(\theta(w))} = \delta \cdot \frac{1}{\alpha}
= \frac{\delta}{\alpha}.
\end{equation}
The result is income-level invariant: the Pareto constant inverse hazard ensures that the
fiscal cost of the Pigouvian correction is proportional to $1/\alpha$ everywhere in the
distribution. Note also the dual role of the Pareto ratio $\alpha/(\alpha-1)$:
equation~\eqref{eq:pareto_hazard} gives $[1-F]/[wf] = 1/\alpha$ (the inverse hazard), while
equation~\eqref{eq:signaling_premium} gives $[w - \underline{\theta}]/w = 1/\alpha$ (the
signaling premium). Both expressions flow from the same Pareto moment identity
$\mathbb{E}[\theta \mid \theta > \underline{\theta}] = \underline{\theta}\cdot\alpha/(\alpha-1)$,
confirmed by the integral in Step~1(b). $\square$

\medskip
\noindent\textit{Note on the redistributive component.} Proposition~\ref{prop:taup} isolates the
Pigouvian component $\tau_p$. The full optimal tax is $t^* = t_{\text{equity}} - \tau_p$, where
$t_{\text{equity}}$ is determined by the standard Mirrlees--Saez formula
\citep{mirrlees1971exploration, saez2001using} applied to redistributive preferences. The two
components are separable under the utilitarian social welfare function: adding the informational
externality shifts the optimum by exactly $\tau_p = \delta/\alpha$ regardless of the redistributive
optimum. \citet{sztutman2024} derives $t_{\text{equity}}$ in full; the contribution of this
paper is to estimate the sufficient statistics $(\hat\delta, \hat\alpha)$ that pin down $\tau_p$.

\bigskip

\subsection*{Derivation of Corollary~\ref{prop:gamma_arc}}

\noindent \textit{Statement.} If $\delta(\text{exp})$ is strictly decreasing in potential
experience, then $\gamma(\text{exp})$ is also strictly decreasing in potential experience.

\noindent \textit{Derivation.} From equation~(\ref{eq:gamma_structural}):
\[
\gamma(\text{exp}) = \frac{\delta(\text{exp})}{1 - \delta(\text{exp}) + \varepsilon^w}.
\]
Let $D \equiv 1 - \delta + \varepsilon^w > 0$ (denominator, positive since $\varepsilon^w > -1$
and $\delta < 1$). Differentiating with respect to experience:
\begin{align*}
\frac{\partial \gamma}{\partial \text{exp}}
&= \frac{\delta'(\text{exp}) \cdot D - \delta(\text{exp}) \cdot (-\delta'(\text{exp}))}{D^2} \\
&= \frac{\delta'(\text{exp}) \cdot (D + \delta(\text{exp}))}{D^2} \\
&= \frac{\delta'(\text{exp}) \cdot (1 + \varepsilon^w)}{D^2}.
\end{align*}
Since $1 + \varepsilon^w > 0$ and $D^2 > 0$, we have
$\text{sign}(\partial \gamma / \partial \text{exp}) = \text{sign}(\delta'(\text{exp}))$.
If $\delta'(\text{exp}) < 0$, then $\gamma'(\text{exp}) < 0$. $\square$

\noindent \textit{Remark.} The corollary does not require Assumption~\ref{ass:pareto}: it
holds under any type distribution, as long as $\varepsilon^w > -1$. The Pareto assumption
is required only for the closed-form formula $\tau_p = \delta/\alpha$ in
Proposition~\ref{prop:taup}; the monotone-$\gamma$ prediction of Corollary~\ref{prop:gamma_arc}
is more general.

%==============================================================================
\newpage
\section{Occupation-Level $\hat{\gamma}$ Estimates}
\label{app:occ_gamma}

%------ Table A1: gamma by 1-digit occupation ------
\begin{table}[htbp]
\centering
\caption{Wage-Cumulative-Hours Elasticity $\hat{\gamma}$ by 1-Digit Occupation}
\label{tab:gamma_occ}
\begin{threeparttable}
\small
\setlength{\tabcolsep}{6pt}
\begin{tabular}{llcccl}
\toprule
Occ.\ Code & Occupation & $\hat{\gamma}_{\text{FE}}$ & SE & $N$ & Theory Group \\
\midrule
3 & Sales                   & 1.235 & 0.107 &  4,955 & Medium \\
4 & Clerical                & 1.077 & 0.042 & 21,653 & Medium \\
1 & Professional/Technical  & 1.050 & 0.062 & 13,392 & High   \\
9 & Service Workers         & 0.968 & 0.035 & 18,628 & Medium \\
7 & Laborers                & 0.937 & 0.058 &  7,447 & Low    \\
6 & Operatives              & 0.878 & 0.042 & 14,424 & Medium \\
2 & Managers/Admin          & 0.867 & 0.087 &  7,874 & High   \\
5 & Craftsmen               & 0.716 & 0.056 & 10,809 & Medium \\
8 & Farm Workers            & 0.470 & 0.174 &  1,474 & Low    \\
\midrule
\multicolumn{3}{l}{\textit{Group averages (from Table~\ref{tab:sector}):}} \\
\quad High (occ 1--2) & & 1.037 & 0.027 & 21,266 & \\
\quad Medium (occ 3--6, 9) & & 0.953 & 0.028 & 70,659 & \\
\quad Low (occ 7--8) & & 0.900 & 0.036 &  8,921 & \\
\bottomrule
\end{tabular}
\begin{tablenotes}
\small
\item \textit{Notes:} FE estimates from equation~(\ref{eq:fe_main}) estimated separately
within each occupation group, NLSY79 1979--1993 annual period only. Sorted by $\hat{\gamma}$
(descending). Theory Group assignments: High signaling = professional/managerial occupations
where output quality is difficult for employers to observe; Low signaling = farm and
laboring occupations where output is directly observable; Medium = remaining occupations.
The Sales occupation has the highest $\hat{\gamma}$, consistent with severe information
asymmetry between salespeople and employers (client-facing output is unobservable to
current employer). Farm Workers have the lowest $\hat{\gamma}$, consistent with directly
observable agricultural output and minimal employer uncertainty. All estimates significant
at $p < 0.01$; standard errors clustered by person.
\end{tablenotes}
\end{threeparttable}
\end{table}

%==============================================================================
\newpage
\section*{Tables and Figures}

%------ Table 1: Descriptive Statistics ------
\begin{table}[htbp]
\centering
\caption{Descriptive Statistics}
\label{tab:desc}
\begin{threeparttable}
\small
\setlength{\tabcolsep}{8pt}
\begin{tabular}{lcccc}
\toprule
Variable & Mean & Std.\ Dev. & Median & $N$ \\
\midrule
Primary wages (1984 \$) & 27,695 & 39,322 & 17,000 & 189,671 \\
Potential experience (years) & 14.0 & 11.7 & 12.0 & 189,671 \\
Highest grade completed & 13.6 & 3.3 & 13.0 & 189,671 \\
Age & 33.4 & 11.9 & 32.0 & 189,671 \\
Cumulative hours worked & 26,344 & 23,627 & 18,544 & 189,671 \\
\midrule
\multicolumn{5}{l}{\textit{Estimation sample restrictions:}} \\
Structural $\gamma$ sample (age 18--65, exp 0--40, pwages $> 0$) & \multicolumn{3}{c}{} & 179,209 \\
Annual ETI sample (\$5K floor, $|\Delta \log z| \leq \log 5$) & \multicolumn{3}{c}{} & 53,632 \\
Biennial ETI sample (\$10K floor, $|\Delta \log z| \leq \log 5$) & \multicolumn{3}{c}{} & 47,201 \\
\bottomrule
\end{tabular}
\begin{tablenotes}
\small
\item \textit{Notes:} Full analytical population is defined as age 18--65 with positive primary
wages, NLSY79 survey rounds 1979--2023. Wages deflated to 1984 dollars using CPI-U.
Cumulative hours accumulated from labor market entry; biennial gaps interpolated linearly
within person.
\end{tablenotes}
\end{threeparttable}
\end{table}

%------ Table 2: Career-Stage gamma ------
\begin{table}[htbp]
\centering
\caption{Wage-Cumulative-Hours Elasticity $\hat{\gamma}$ by Career Stage}
\label{tab:gamma}
\begin{threeparttable}
\small
\setlength{\tabcolsep}{5pt}
\begin{tabular}{lccccc}
\toprule
 & & \multicolumn{2}{c}{FE Estimator} & OLS & \\
\cmidrule(lr){3-4}
Career Stage & $N$ & $\hat{\gamma}_{\text{FE}}$ & (SE) & $\hat{\gamma}_{\text{OLS}}$ & Selection \\
\midrule
Q1: exp 0--8 yr    & 39,520  & 1.101 & (0.021) & 0.991 & $-$0.110 \\
Q2: exp 8--13 yr   & 38,814  & 1.049 & (0.038) & 0.964 & $-$0.085 \\
Q3: exp 13--17 yr  & 32,892  & 1.067 & (0.063) & 1.030 & $-$0.037 \\
Q4: exp 17--25 yr  & 34,383  & 0.927 & (0.057) & 0.972 & $+$0.044 \\
Q5: exp $> 25$ yr  & 33,600  & 1.115$^{\dagger}$ & (0.110) & 1.086 & $-$0.029 \\
\midrule
Full career        & 179,209 & 0.937 & (0.011) & 1.022 & $+$0.086 \\
\midrule
\multicolumn{6}{l}{\textit{Polynomial specification: $\hat{\gamma}(\text{exp}) = \hat{\beta}_0 + \hat{\beta}_1\cdot\text{exp} + \hat{\beta}_2\cdot\text{exp}^2$}} \\
$\hat{\beta}_0$ (entry) & & 0.974 & (0.011) & & \\
$\hat{\beta}_1$ (linear decline) & & $-$0.0134 & (0.0015) & & \\
$\hat{\beta}_2$ (quadratic) & & $+$0.0003 & (0.00004) & & \\
\multicolumn{2}{l}{Joint test $\beta_1 = \beta_2 = 0$:} & \multicolumn{4}{l}{$F(2, 11{,}524) = 43.20$\quad ($p < 0.001$)} \\
\bottomrule
\end{tabular}
\begin{tablenotes}
\small
\item \textit{Notes:} FE estimates from equation~(\ref{eq:fe_main}) with person fixed effects,
quadratic potential experience, and year fixed effects. Standard errors clustered by person.
Selection $= \hat{\gamma}_{\text{OLS}} - \hat{\gamma}_{\text{FE}}$; positive values indicate
positive selection on ability (high-ability workers supply more hours). $^{\dagger}$Q5 is
non-identified: within-$R^2 = 2.6\%$; 95\% CI $= [0.90, 1.33]$ contains all Q1--Q4 estimates;
survivor selection gap $= +0.567$ log-wage points (Section~\ref{sec:robustness}). Q5 is
excluded from the structural career arc.
\end{tablenotes}
\end{threeparttable}
\end{table}

%------ Table 3: ETI ------
\begin{table}[htbp]
\centering
\caption{Elasticity of Taxable Income: Two-Period IV Estimates}
\label{tab:eti}
\begin{threeparttable}
\small
\setlength{\tabcolsep}{4pt}
\begin{tabular}{l p{3.5cm} c c c c c}
\toprule
Period & Specification & $N$ & $\hat{\varepsilon}^w$ & SE & $t$-stat & 1st-stage $F$ \\
\midrule
Annual (1978--93) & Primary (\$5K floor) & 53,632 & $-$0.048 & 0.080 & $-$0.60 & 2,617 \\
Annual (1978--93) & Near-worker (\$1K floor) & 59,428 & $-$0.049 & 0.082 & $-$0.60 & 2,500 \\
Annual (1978--93) & Kopczuk lagged control & 1,595 & $+$0.021 & 0.352 & $+$0.06 & 149 \\
\midrule
Biennial (1995--19) & Primary (\$10K floor) & 47,201 & $+$0.297 & 0.121 & $+$2.46 & 1,146 \\
\bottomrule
\end{tabular}
\begin{tablenotes}
\small
\item \textit{Notes:} IV estimates of equation~(\ref{eq:eti}). Instrument is the simulated
change in log net-of-tax rate, holding individual income at its base-period level inflated by
average earnings growth. Annual period uses $k = 3$-year horizon; biennial uses $k = 2$-period
horizon. All specifications include log base-period income as mean-reversion control, year fixed
effects, and controls for marital status transitions. Standard errors clustered by person.
Identifying reforms: ERTA (1981), TRA (1986) [annual]; EGTRRA (2001), JGTRRA (2003),
TCJA (2017) [biennial]. Kopczuk specification additionally controls for one-period lagged
income; sample collapses to N $= 1{,}595$ balanced two-window observations.
\end{tablenotes}
\end{threeparttable}
\end{table}

%------ Table 4: Pigouvian Tax ------
\begin{table}[htbp]
\centering
\caption{Optimal Pigouvian Labor Income Subsidy: $\hat{\tau}_p = \hat{\delta}/\hat{\alpha}$}
\label{tab:taup}
\begin{threeparttable}
\small
\begin{tabular}{p{6cm} c c c c}
\toprule
Specification & $\hat{\delta}$ & $\hat{\alpha}$ & $\hat{\tau}_p$ & 95\% CI \\
\midrule
\textit{Primary estimate} & & & & \\
\quad Biennial ETI ($\hat{\varepsilon}^w = +0.297$) & 0.628 & 3.295 & 19.1\% & $[15.6\%, 22.6\%]$ \\
\midrule
\textit{Lifecycle lower bound} & & & & \\
\quad Annual ETI ($\hat{\varepsilon}^w = -0.048$) & 0.461 & 3.295 & 14.0\% & $[11.7\%, 16.3\%]$ \\
\midrule
\textit{IV lower bound (simultaneity-adjusted)} & & & & \\
\quad $\hat{\gamma}_{\text{IV}} = 0.697$, biennial $\hat{\varepsilon}^w$ & 0.533 & 3.295 & 16.2\% & --- \\
\midrule
\textit{Sensitivity to $\hat{\alpha}$} (biennial $\hat{\delta} = 0.628$) & & & & \\
\quad $\alpha = 3.0$ & 0.628 & 3.0 & 20.9\% & --- \\
\quad $\alpha = 2.5$ & 0.628 & 2.5 & 25.1\% & --- \\
\quad $\alpha = 1.5$ & 0.628 & 1.5 & 41.9\% & --- \\
\midrule
\textit{Comparison} & & & & \\
\quad Sztutman (2024), HRS ages 50+ & $\approx 0.17$ & --- & $\approx 5\%$ & --- \\
\bottomrule
\end{tabular}
\begin{tablenotes}
\small
\item \textit{Notes:} Primary estimate uses $\hat{\gamma}_{\text{FE}} = 0.937$ (full-career
FE), biennial ETI $\hat{\varepsilon}^w = 0.297$, and $\hat{\alpha} = 3.295$ from log-rank
regression ($R^2 = 0.985$). Delta-method CI propagates uncertainty through
$\hat{\delta} = \hat{\gamma}(1 + \hat{\varepsilon}^w)/(1 + \hat{\gamma})$; the biennial ETI
contributes 99.6\% of variance in $\hat{\tau}_p$. Bootstrap CI (B $= 500$, clustered) agrees
within $< 1$ pp. IV lower bound uses $\hat{\gamma}_{\text{IV}} = 0.697$ to address within-year
simultaneity. Sztutman (2024) comparison is a back-of-envelope estimate from the HRS sample.
\end{tablenotes}
\end{threeparttable}
\end{table}

%------ Table 5: Sector Heterogeneity and AP Test ------
\begin{table}[htbp]
\centering
\caption{Sector Heterogeneity: $\hat{\gamma}$ and Altonji--Pierret Slopes by Signaling Group}
\label{tab:sector}
\begin{threeparttable}
\small
\setlength{\tabcolsep}{5pt}
\begin{tabular}{lccccc}
\toprule
 & & \multicolumn{2}{c}{$\hat{\gamma}_{\text{FE}}$} & \multicolumn{2}{c}{AP Slope ($\hat{\phi}_2$)} \\
\cmidrule(lr){3-4}\cmidrule(lr){5-6}
Group & Occupations & Estimate & SE & Full sample & 5+yr tenure \\
\midrule
Low (farm/laborers) & 7--8 & 0.900 & 0.036 & 0.00511$^{**}$ & 0.00432 \\
Medium (clerical/sales/ops/service) & 3--6, 9 & 0.953 & 0.028 & 0.00960$^{***}$ & 0.00499 \\
High (prof/tech, managers) & 1--2 & 1.037 & 0.027 & 0.00666$^{***}$ & 0.00882$^{**}$ \\
\midrule
\multicolumn{2}{l}{High $-$ Low gap:} & 0.137 & 0.045 & & \\
\multicolumn{2}{l}{95\% CI for High $-$ Low:} & \multicolumn{2}{l}{$[0.067, 0.207]$} & & \\
\bottomrule
\end{tabular}
\begin{tablenotes}
\small
\item \textit{Notes:} $\hat{\gamma}$ from equation~(\ref{eq:fe_main}) within occupation group,
NLSY79 1979--1993. AP slope is $\hat{\phi}_2$ from equation~(\ref{eq:ap}): the coefficient
on AFQT $\times$ potential experience. ``Full sample'' uses all workers in each group.
``5+yr tenure'' restricts to workers with at least 5 years of occupation-specific experience.
AP slopes in the ``5+yr tenure'' column are group averages across constituent occupations
(e.g., High group $= $ average of Professional/Tech and Managers/Admin). The full-sample
Low-group spike (0.00511) is driven by Laborers' short-tenure workers; restricting to
5+ year tenure, Laborers collapse to 0.00015 ($t = 0.01$), while Managers/Admin rise to
0.01493 ($t = 2.85$). Group-level averages understate these within-group contrasts.
Standard errors clustered by person. $^{**} p < 0.05$; $^{***} p < 0.01$.
\end{tablenotes}
\end{threeparttable}
\end{table}

%------ Table 6: Delta Lifecycle Arc ------
\begin{table}[htbp]
\centering
\caption{$\hat{\delta}$ Lifecycle Arc: Implied Pigouvian Tax by Career Stage}
\label{tab:delta_arc}
\begin{threeparttable}
\small
\setlength{\tabcolsep}{6pt}
\begin{tabular}{lcccccc}
\toprule
 & & \multicolumn{2}{c}{Annual ETI ($\hat{\varepsilon}^w = -0.048$)} & \multicolumn{2}{c}{Biennial ETI ($\hat{\varepsilon}^w = +0.297$)} \\
\cmidrule(lr){3-4}\cmidrule(lr){5-6}
Exp (yr) & $\hat{\gamma}(\text{exp})$ & $\hat{\delta}$ & $\hat{\tau}_p$ & $\hat{\delta}$ & $\hat{\tau}_p$ \\
\midrule
0  & 0.974 & 0.470 & 14.3\% & 0.640 & 19.4\% \\
5  & 0.915 & 0.455 & 13.8\% & 0.619 & 18.8\% \\
10 & 0.871 & 0.443 & 13.4\% & 0.604 & 18.3\% \\
15 & 0.843 & 0.435 & 13.2\% & 0.593 & 18.0\% \\
20 & 0.830 & 0.432 & 13.1\% & 0.588 & 17.8\% \\
25 & 0.833 & 0.433 & 13.1\% & 0.589 & 17.9\% \\
30 & 0.851 & 0.438 & 13.3\% & 0.596 & 18.1\% \\
35 & 0.886 & 0.447 & 13.6\% & 0.609 & 18.5\% \\
40 & 0.935 & 0.460 & 14.0\% & 0.627 & 19.0\% \\
\midrule
\multicolumn{2}{l}{Within-regime range:} & \multicolumn{2}{c}{1.3 pp} & \multicolumn{2}{c}{1.6 pp} \\
\bottomrule
\end{tabular}
\begin{tablenotes}
\small
\item \textit{Notes:} $\hat{\gamma}(\text{exp})$ from polynomial specification
(\ref{eq:fe_poly}). $\hat{\delta}(\text{exp})$ from equation~(\ref{eq:delta}) at fixed
$\hat{\varepsilon}^w$. $\hat{\tau}_p(\text{exp}) = \hat{\delta}(\text{exp})/\hat{\alpha}$
with $\hat{\alpha} = 3.295$. Entries at exp $\geq 25$ include the slight quadratic upturn in
$\hat{\gamma}$ and partly reflect the survivor selection documented in
Section~\ref{sec:robustness}; the primary structural arc is exp 0--25.
\end{tablenotes}
\end{threeparttable}
\end{table}

%------ Figure 1: gamma arc ------
\begin{figure}[htbp]
\centering
\caption{Career Arc of the Wage-Cumulative-Hours Elasticity $\hat{\gamma}(\text{exp})$}
\label{fig:gamma}
\includegraphics[width=0.85\textwidth]{output/Phase1_Fig3_PolynomialGamma.png}
\begin{minipage}{0.85\textwidth}
\small
\textit{Notes:} Estimated from polynomial specification (\ref{eq:fe_poly}) on the structural
sample ($N = 179{,}209$). The solid line is $\hat{\gamma}(\text{exp}) = 0.974 - 0.0134 \cdot
\text{exp} + 0.0003 \cdot \text{exp}^2$. Shaded region is the pointwise 95\% confidence band.
Career minimum at exp$^* = 21.6$ years, where $\hat{\gamma} = 0.829$. The quadratic upturn
after exp $= 25$ reflects survivor selection at late career (Section~\ref{sec:robustness})
and is not part of the main structural arc.
\end{minipage}
\end{figure}

\end{document}
